% ============================================================
% HSBI --- Quantitative Research Methods
% FINAL Mock Examination (Lectures 1--12, comprehensive)
% Focus: ISLP Chapters 7, 8, 10, 13 (~60%) + Chapters 2--6 (~40%)
% Prof. Dr. Christoph Weisser --- Summer Semester 2026
% ------------------------------------------------------------
% SINGLE SOURCE with solutions toggle:
%   Exam only:
%     pdflatex -jobname=Final_Mock_Exam final_mock_exam.tex
%   With solutions:
%     pdflatex -jobname=Final_Mock_Exam_Solutions "\def\withsolutions{1}% ============================================================
% HSBI --- Quantitative Research Methods
% FINAL Mock Examination (Lectures 1--12, comprehensive)
% Focus: ISLP Chapters 7, 8, 10, 13 (~60%) + Chapters 2--6 (~40%)
% Prof. Dr. Christoph Weisser --- Summer Semester 2026
% ------------------------------------------------------------
% SINGLE SOURCE with solutions toggle:
%   Exam only:
%     pdflatex -jobname=Final_Mock_Exam final_mock_exam.tex
%   With solutions:
%     pdflatex -jobname=Final_Mock_Exam_Solutions "\def\withsolutions{1}% ============================================================
% HSBI --- Quantitative Research Methods
% FINAL Mock Examination (Lectures 1--12, comprehensive)
% Focus: ISLP Chapters 7, 8, 10, 13 (~60%) + Chapters 2--6 (~40%)
% Prof. Dr. Christoph Weisser --- Summer Semester 2026
% ------------------------------------------------------------
% SINGLE SOURCE with solutions toggle:
%   Exam only:
%     pdflatex -jobname=Final_Mock_Exam final_mock_exam.tex
%   With solutions:
%     pdflatex -jobname=Final_Mock_Exam_Solutions "\def\withsolutions{1}% ============================================================
% HSBI --- Quantitative Research Methods
% FINAL Mock Examination (Lectures 1--12, comprehensive)
% Focus: ISLP Chapters 7, 8, 10, 13 (~60%) + Chapters 2--6 (~40%)
% Prof. Dr. Christoph Weisser --- Summer Semester 2026
% ------------------------------------------------------------
% SINGLE SOURCE with solutions toggle:
%   Exam only:
%     pdflatex -jobname=Final_Mock_Exam final_mock_exam.tex
%   With solutions:
%     pdflatex -jobname=Final_Mock_Exam_Solutions "\def\withsolutions{1}\input{final_mock_exam.tex}"
% ============================================================
\documentclass[11pt,a4paper]{article}
\usepackage[utf8]{inputenc}
\usepackage[T1]{fontenc}
\usepackage[english]{babel}
\usepackage[margin=2.2cm]{geometry}
\usepackage{amsmath,amssymb}
\usepackage{booktabs,array}
\usepackage{enumitem}
\usepackage{xcolor}
\usepackage{tcolorbox}
\tcbuselibrary{skins,breakable}

\setlength{\parindent}{0pt}

% Teal solution box (only shown when \withsolutions is defined)
\newtcolorbox{solutionbox}{enhanced,breakable,colback=teal!5,colframe=teal!55!black,
  boxrule=0pt,leftrule=2.5pt,arc=2pt,fonttitle=\bfseries,title=Solution,
  left=2mm,right=2mm,top=1mm,bottom=1.5mm}

\newcommand{\problem}[2]{\par\vspace{7mm}\noindent{\large\textbf{Problem #1}}%
  \hfill{\large\textbf{(#2 points)}}\par\nopagebreak\vspace{0.6mm}%
  \noindent\rule{\textwidth}{0.6pt}\par\nopagebreak\vspace{2.5mm}}
\newcommand{\pp}[1]{\textbf{[#1~P]}\enspace}

\begin{document}

% ------------------------------------------------------------
% Header block
% ------------------------------------------------------------
\begin{center}
  {\Large\textbf{HSBI --- Bielefeld University of Applied Sciences and Arts}}\\[1.2mm]
  {\large Quantitative Research Methods}\\[1mm]
  {\normalsize Summer Semester 2026 \quad$\cdot$\quad Examiner: Prof.\ Dr.\ Christoph Weisser}\\[4.5mm]
  {\LARGE\textbf{Final Mock Examination (Lectures 1--12)}}\\[1.2mm]
  {\normalsize Comprehensive coverage of ISLP Chapters 2--8, 10 and 13,\\
   with emphasis on Chapters 7 (splines/GAMs), 8 (trees), 10 (deep learning) and 13 (multiple testing)}%
  \ifdefined\withsolutions\\[2.2mm]{\large\textcolor{teal!55!black}{\textbf{--- Version with detailed solutions ---}}}\fi
\end{center}

\vspace{2.5mm}
\begin{center}
\begin{tabular}{@{}ll@{}}
  \textbf{Duration:} & \textbf{120 minutes}\\
  \textbf{Total credit:} & \textbf{120 points}\\
  \textbf{Allowed aids:} & non-programmable calculator; one handwritten A4 sheet of notes (both sides)\\
\end{tabular}
\end{center}

\vspace{3mm}
Name: \rule[-1mm]{0.38\textwidth}{0.4pt}\hfill Matriculation no.: \rule[-1mm]{0.24\textwidth}{0.4pt}

\vspace{4.5mm}
\begin{center}
\renewcommand{\arraystretch}{1.35}
\begin{tabular}{|l|c|c|c|c|c|c|c|c|}
\hline
\textbf{Problem} & \textbf{P1} & \textbf{P2} & \textbf{P3} & \textbf{P4} & \textbf{P5} & \textbf{P6} & \textbf{P7} & \textbf{Total}\\\hline
Maximum points & 16 & 8 & 20 & 16 & 20 & 20 & 20 & \textbf{120}\\\hline
Achieved points & \hspace{8mm} & \hspace{8mm} & \hspace{8mm} & \hspace{8mm} & \hspace{8mm} & \hspace{8mm} & \hspace{8mm} & \hspace{10mm}\\\hline
\end{tabular}
\end{center}

\vspace{2.5mm}
\textbf{Instructions}
\begin{itemize}[nosep,leftmargin=5.5mm]
  \item Justify all answers. Numerical results without a derivation or a brief reasoning receive no credit.
  \item Round final numerical results to 3 decimal places unless stated otherwise. Small deviations caused by carrying rounded intermediate results are accepted.
  \item All numerical constants you need (powers of $e$, logarithms, $(0.95)^{20}$, $t$-quantiles) are provided in the problems.
  \item The number of points per problem roughly equals the recommended working time in minutes.
\end{itemize}
\vspace{1mm}\noindent\rule{\textwidth}{0.6pt}

% ============================================================
\problem{1: Moving beyond linearity --- splines}{16}

\begin{enumerate}[label=(\alph*),leftmargin=7mm,itemsep=2.5mm]
  \item \pp{4} A \emph{cubic spline} is a piecewise cubic polynomial that is continuous and has
    continuous first and second derivatives at each knot. How many degrees of freedom (free
    parameters, counting the intercept) does a cubic spline with $K$ interior knots use?
    Justify your answer by counting the parameters of the polynomial pieces and the constraints
    at the knots, and evaluate the result for $K=4$.
  \item \pp{4} How many degrees of freedom does a \emph{natural} cubic spline with the same $K$
    knots use (again evaluate for $K=4$)? Which additional constraint does a natural spline
    impose, in which region of the predictor space does it act, and why is this constraint
    useful in practice?
  \item \pp{4} A \emph{smoothing spline} is the function $g$ that minimises
    \[
      \sum_{i=1}^{n}\bigl(y_i-g(x_i)\bigr)^2+\lambda\int g''(t)^2\,dt .
    \]
    Explain the role of the two terms, and describe the behaviour of the solution $\hat{g}$ in
    the two limiting cases $\lambda\to 0$ and $\lambda\to\infty$.
  \item \pp{4} Write down the \emph{truncated-power basis} representation of a cubic spline with
    a single knot at $\xi$, define the truncated-power basis function explicitly, and explain in
    one sentence why adding this basis function does not destroy the smoothness of the fit at
    $\xi$.
\end{enumerate}

\ifdefined\withsolutions
\begin{solutionbox}
\textbf{(a) Degrees of freedom of a cubic spline.}
$K$ interior knots split the range of $X$ into $K+1$ intervals, each carrying a cubic
polynomial with $4$ coefficients: $4(K+1)=4K+4$ parameters. At each knot, continuity of
$g$, $g'$ and $g''$ imposes $3$ constraints, i.e.\ $3K$ constraints in total. Hence
\[
  \mathrm{df}=4(K+1)-3K=K+4 .
\]
For $K=4$: $\mathrm{df}=4+4=\mathbf{8}$ free parameters (equivalently: intercept, $x$, $x^2$,
$x^3$ plus one truncated-power basis function per knot $=4+K$).

\medskip
\textbf{(b) Natural cubic spline.}
A natural spline additionally requires the function to be \emph{linear beyond the boundary
knots}, i.e.\ in the two outer regions $X<\xi_1$ and $X>\xi_K$. Forcing the quadratic and cubic
terms to vanish there removes $2$ parameters per boundary region, i.e.\ $4$ in total:
\[
  \mathrm{df}=(K+4)-4=K,\qquad K=4:\ \mathrm{df}=\mathbf{4}.
\]
Usefulness: polynomials and unconstrained splines are notoriously erratic near the boundaries,
where data are sparse; the linearity constraint stabilises the fit there
(\emph{lower variance} of predictions at and beyond the boundary), at the price of a small
amount of additional bias.

\medskip
\textbf{(c) Smoothing spline.}
The first term (RSS) rewards fidelity to the data; the second term penalises roughness, since
$g''$ measures the curvature (wiggliness) of $g$, and $\lambda\ge 0$ is the tuning parameter
that trades the two off.
\emph{$\lambda\to 0$:} the penalty disappears; $\hat{g}$ can be made arbitrarily wiggly and
\emph{interpolates} the training observations exactly (zero training RSS, severe overfitting;
effective df $\to n$).
\emph{$\lambda\to\infty$:} any curvature is infinitely expensive, so $g''\equiv 0$ is enforced;
$\hat{g}$ becomes a straight line and equals the \emph{least-squares regression line}
(effective df $\to 2$). Between the extremes, $\lambda$ (chosen e.g.\ by LOOCV) controls the
effective degrees of freedom, which decrease smoothly from $n$ to $2$.

\medskip
\textbf{(d) Truncated-power basis with one knot $\xi$.}
\[
  f(x)=\beta_0+\beta_1x+\beta_2x^2+\beta_3x^3+\beta_4\,(x-\xi)^3_+,
  \qquad
  (x-\xi)^3_+=\begin{cases}(x-\xi)^3, & x>\xi\\[0.5mm] 0, & x\le\xi\end{cases}
\]
The function $(x-\xi)^3_+$ and its first and second derivatives are all $0$ at $x=\xi$, so
adding it changes only the third derivative at the knot; the fit remains continuous with
continuous first and second derivatives --- exactly the defining smoothness of a cubic spline.
\end{solutionbox}
\fi

% ============================================================
\problem{2: Generalised additive models}{8}

\begin{enumerate}[label=(\alph*),leftmargin=7mm,itemsep=2.5mm]
  \item \pp{3} Write down the general form of a \emph{generalised additive model} (GAM) for a
    quantitative response $Y$ and predictors $X_1,\dots,X_p$, and explain in one sentence how it
    extends multiple linear regression.
  \item \pp{3} State one important \emph{advantage} of GAMs that concerns the interpretation of
    the fitted functions $f_j$, and the main structural \emph{limitation} of GAMs --- including
    how this limitation can be repaired.
  \item \pp{2} Describe in one or two sentences how the \emph{backfitting} algorithm fits a GAM.
\end{enumerate}

\ifdefined\withsolutions
\begin{solutionbox}
\textbf{(a) Model form.}
\[
  Y=\beta_0+f_1(X_1)+f_2(X_2)+\dots+f_p(X_p)+\varepsilon .
\]
Multiple linear regression is the special case $f_j(X_j)=\beta_jX_j$; a GAM replaces each linear
component by a smooth non-linear function $f_j$ (e.g.\ a natural or smoothing spline, or local
regression), while keeping the \emph{additive} structure.

\medskip
\textbf{(b) Advantage and limitation.}
\emph{Advantage (interpretability):} because the model is additive, the contribution of each
$X_j$ to $Y$ is captured by its own function $f_j$, which can be examined and plotted
individually while the other variables are held fixed --- non-linear effects remain as easy to
read as in a coefficient table, one panel per predictor.
\emph{Limitation:} additivity excludes \emph{interactions}: the effect of $X_j$ cannot depend
on the value of $X_k$. Repair: add interaction terms manually, e.g.\ a product term
$X_j\times X_k$ or a low-dimensional smooth surface $f_{jk}(X_j,X_k)$ as an extra component.

\medskip
\textbf{(c) Backfitting.}
Cycle through the predictors and repeatedly re-fit each $f_j$ (by a smoother, e.g.\ a spline)
to the \emph{partial residuals} $r_i=y_i-\beta_0-\sum_{k\ne j}f_k(x_{ik})$, holding all other
fitted functions fixed; iterate until the $f_j$ stop changing.
\end{solutionbox}
\fi

% ============================================================
\problem{3: Regression and classification trees by hand}{20}

An online platform records for $n=6$ products the advertising budget $x$ (in 1{,}000 EUR) and
the weekly sales $y$ (in 1{,}000 units):

\begin{center}
\begin{tabular}{lcccccc}
\toprule
Product $i$ & 1 & 2 & 3 & 4 & 5 & 6\\
\midrule
Budget $x_i$ & 1 & 2 & 3 & 4 & 5 & 6\\
Sales $y_i$ & 2 & 4 & 3 & 3 & 8 & 10\\
\bottomrule
\end{tabular}
\end{center}

A regression tree with a \emph{single} split at cutpoint $s$ partitions the data into
$R_1=\{i: x_i<s\}$ and $R_2=\{i: x_i\ge s\}$ and predicts the region mean
$\hat{y}_{R_m}$ in each region.

\begin{enumerate}[label=(\alph*),leftmargin=7mm,itemsep=2.5mm]
  \item \pp{8} For each of the two candidate cutpoints $s=2.5$ and $s=4.5$, determine the two
    regions, the region means and the resulting total residual sum of squares
    $\operatorname{RSS}(s)=\sum_{m=1}^{2}\sum_{i\in R_m}(y_i-\hat{y}_{R_m})^2$.
    Which cutpoint does recursive binary splitting choose, and why?
  \item \pp{2} Using the chosen split, predict the weekly sales for a new product with
    advertising budget $x=5$.
  \item \pp{6} A \emph{classification} tree for a different task has a terminal node containing
    8 training observations: 6 of class ``Yes'' and 2 of class ``No''. Compute the class
    proportions, the Gini index $G=\sum_{k}\hat{p}_{k}(1-\hat{p}_{k})$ and the entropy
    $D=-\sum_{k}\hat{p}_{k}\ln\hat{p}_{k}$ of this node
    \emph{(use $\ln 0.75=-0.2877$ and $\ln 0.25=-1.3863$)}.
    Which class does the node predict?
  \item \pp{4} In \emph{cost-complexity pruning} one minimises, for a subtree
    $T\subseteq T_0$ of the large tree $T_0$,
    \[
      \sum_{m=1}^{|T|}\ \sum_{i:\,x_i\in R_m}\bigl(y_i-\hat{y}_{R_m}\bigr)^2+\alpha\,|T| .
    \]
    Explain the role of the tuning parameter $\alpha$: what happens for $\alpha=0$, what happens
    as $\alpha$ increases, and how is $\alpha$ chosen in practice?
\end{enumerate}

\ifdefined\withsolutions
\begin{solutionbox}
\textbf{(a) Evaluating the two cutpoints.}
Root node for reference: $\bar{y}=30/6=5$.

\emph{Cutpoint $s=2.5$:} $R_1=\{1,2\}$ with $y$-values $\{2,4\}$, mean $\hat{y}_{R_1}=3$;
$R_2=\{3,4,5,6\}$ with $y$-values $\{3,3,8,10\}$, mean $\hat{y}_{R_2}=24/4=6$.
\[
  \operatorname{RSS}_{R_1}=(2-3)^2+(4-3)^2=2,\quad
  \operatorname{RSS}_{R_2}=(3-6)^2+(3-6)^2+(8-6)^2+(10-6)^2=9+9+4+16=38,
\]
so $\operatorname{RSS}(2.5)=2+38=\mathbf{40}$.

\emph{Cutpoint $s=4.5$:} $R_1=\{1,2,3,4\}$ with $y$-values $\{2,4,3,3\}$, mean
$\hat{y}_{R_1}=12/4=3$; $R_2=\{5,6\}$ with $y$-values $\{8,10\}$, mean $\hat{y}_{R_2}=9$.
\[
  \operatorname{RSS}_{R_1}=1+1+0+0=2,\qquad
  \operatorname{RSS}_{R_2}=(8-9)^2+(10-9)^2=2,
\]
so $\operatorname{RSS}(4.5)=2+2=\mathbf{4}$.

Since $4<40$, recursive binary splitting is \emph{greedy} and picks the cutpoint with the
smallest RSS: $\mathbf{s=4.5}$. (It separates the four low-sales products, mean $3$, from the
two high-sales products, mean $9$.)

\medskip
\textbf{(b) Prediction for $x=5$.}
$x=5\ge 4.5$, so the new product falls into $R_2$ and the tree predicts
$\hat{y}=\hat{y}_{R_2}=\mathbf{9}$ (i.e.\ $9{,}000$ units per week).

\medskip
\textbf{(c) Node impurity.}
$\hat{p}_{\text{Yes}}=6/8=0.75$, $\hat{p}_{\text{No}}=2/8=0.25$.
\[
  G=0.75\cdot 0.25+0.25\cdot 0.75=0.1875+0.1875=\mathbf{0.375}.
\]
\[
  D=-\bigl(0.75\ln 0.75+0.25\ln 0.25\bigr)
   =-\bigl(0.75\cdot(-0.2877)+0.25\cdot(-1.3863)\bigr)
   =0.2158+0.3466=\mathbf{0.562}.
\]
The node predicts the majority class, \textbf{Yes} (with estimated probability $0.75$).
Both $G$ and $D$ are small when the node is pure; here the node is moderately impure.

\medskip
\textbf{(d) Role of $\alpha$.}
$\alpha$ is the price per terminal node: it trades off the fit of the subtree (first term)
against its complexity $|T|$.
For $\alpha=0$ the criterion equals the training RSS and the full tree $T_0$ is optimal.
As $\alpha$ increases, additional leaves must ``pay for themselves''; branches whose RSS
reduction is smaller than $\alpha$ are pruned, so the optimal subtree becomes smaller ---
one obtains a \emph{nested sequence} of subtrees as $\alpha$ grows (weakest-link pruning).
In practice $\alpha$ is selected by \emph{cross-validation} (choose the $\alpha$ with the
smallest CV error, then refit on all data), which balances bias against variance instead of
optimising training fit.
\end{solutionbox}
\fi

% ============================================================
\problem{4: Ensembles --- bagging, random forests, boosting}{16}

\begin{enumerate}[label=(\alph*),leftmargin=7mm,itemsep=2.5mm]
  \item \pp{6} Let $\hat{f}_1(x),\dots,\hat{f}_B(x)$ be $B$ identically distributed (but not
    independent) tree predictions at a fixed point $x$, each with variance $\sigma^2$ and with
    pairwise correlation $\rho$. The variance of the bagged average is
    \[
      \operatorname{Var}\Bigl(\tfrac{1}{B}\sum_{b=1}^{B}\hat{f}_b(x)\Bigr)
      =\rho\,\sigma^{2}+\frac{1-\rho}{B}\,\sigma^{2}.
    \]
    (i) Explain briefly why this formula shows that bagging reduces variance, and which term
    limits the reduction. (ii) Evaluate the variance for $B=100$, $\rho=0.6$ and $\sigma^2=4$,
    and compare it with the variance of a single tree. (iii) What value does the variance
    approach as $B\to\infty$?
  \item \pp{4} How do \emph{random forests} modify bagging at each split, and what is the
    recommended subset size $m$ for classification as a function of $p$ (evaluate for $p=25$)?
    Explain, with reference to the formula in (a), why this modification improves on bagging.
  \item \pp{4} Name the three tuning parameters of \emph{boosting} for regression trees and
    state in one sentence each what they control. What is the ``slow learning'' idea behind
    boosting?
  \item \pp{2} What does the \emph{out-of-bag} (OOB) error of a bagged model estimate, and why
    is it available essentially for free?
\end{enumerate}

\ifdefined\withsolutions
\begin{solutionbox}
\textbf{(a) Variance of the bagged average.}
(i) The second term $\frac{1-\rho}{B}\sigma^2$ is driven to $0$ by averaging over many trees
($B$ large), so averaging removes the ``independent part'' of the trees' variability ---
this is the variance reduction of bagging. The first term $\rho\sigma^2$ does \emph{not}
depend on $B$: because the bootstrap trees are grown on overlapping samples of the same data
(and share the same strong predictors), they are positively correlated, and this correlated
part of the variance cannot be averaged away. It limits the reduction.
(ii) For $B=100$, $\rho=0.6$, $\sigma^2=4$:
\[
  \operatorname{Var}=0.6\cdot 4+\frac{1-0.6}{100}\cdot 4
  =2.4+\frac{0.4\cdot 4}{100}=2.4+0.016=\mathbf{2.416},
\]
compared with $\sigma^2=4$ for a single tree --- a reduction of about $40\%$.
(iii) As $B\to\infty$ the variance approaches the floor $\rho\sigma^2=0.6\cdot 4=\mathbf{2.4}$.
(Increasing $B$ further does not overfit; it only stabilises the average.)

\medskip
\textbf{(b) Random forests.}
At each split, only a \emph{random subset} of $m$ of the $p$ predictors is considered as
splitting candidates; a fresh subset is drawn at every split. Recommended for classification:
$m\approx\sqrt{p}$, so for $p=25$: $m=\sqrt{25}=\mathbf{5}$.
Why it helps: with bagging, a dominant predictor is used in the top split of almost every tree,
so the trees look alike and $\rho$ is high. Excluding it from $\approx (p-m)/p$ of the split
searches \emph{decorrelates} the trees, i.e.\ lowers $\rho$ --- and by (a) this lowers the
irreducible floor $\rho\sigma^2$ of the ensemble variance, which mere averaging cannot touch.

\medskip
\textbf{(c) Boosting.}
\begin{itemize}[nosep,leftmargin=5mm]
  \item \emph{Number of trees $B$:} how many small trees are added sequentially; unlike
    bagging, boosting \emph{can} overfit if $B$ is too large, so $B$ is chosen by
    cross-validation.
  \item \emph{Shrinkage $\lambda$} (learning rate, e.g.\ $0.01$ or $0.001$): scales the
    contribution of each new tree and thus the speed of learning; smaller $\lambda$ usually
    needs larger $B$.
  \item \emph{Number of splits $d$ per tree} (interaction depth): the complexity of each tree
    and the interaction order of the model; often even $d=1$ (stumps, additive model) works
    well.
\end{itemize}
\emph{Slow learning:} each tree is fitted to the current \emph{residuals}, and only a small
fraction $\lambda$ of it is added to the model; the model thus improves slowly and exactly in
the regions where it still performs poorly, instead of fitting the data hard in one step ---
which tends to generalise better.

\medskip
\textbf{(d) OOB error.}
Each bootstrap sample leaves out on average about $1/3$ of the observations
(``out-of-bag''). Predicting each observation only with the trees that did \emph{not} use it,
and aggregating these predictions, yields the OOB error --- a valid estimate of the
\emph{test error} of the bagged model. It is free because it reuses the bootstrap side
product; no separate validation set or cross-validation loop is required.
\end{solutionbox}
\fi

% ============================================================
\problem{5: Neural networks by hand}{20}

Consider a tiny feed-forward network for a quantitative response with input
$x=(x_1,x_2)=(1,2)$, one hidden layer with two ReLU units
($g(z)=\max\{0,z\}$) and a linear output unit:
\[
  A_1=g\bigl(-1+1\,x_1+2\,x_2\bigr),\qquad
  A_2=g\bigl(-3-2\,x_1+1\,x_2\bigr),\qquad
  f(x)=1+2\,A_1+4\,A_2 .
\]

\begin{enumerate}[label=(\alph*),leftmargin=7mm,itemsep=2.5mm]
  \item \pp{8} Perform the forward pass step by step: compute both pre-activations, both
    activations $A_1,A_2$ and the network output $f(x)$. Comment briefly on what the ReLU did
    to the second unit and why such non-linear activations are essential.
  \item \pp{6} A classification network with $K=3$ classes produces the logits (output-layer
    pre-activations) $z=(z_1,z_2,z_3)=(1.0,\ 0.0,\ -1.0)$ for one observation whose true class
    is class~1. Compute all three softmax probabilities and the cross-entropy loss
    $-\ln \hat{p}_{\text{true}}$.
    \emph{(Use $e^{1}=2.718$, $e^{0}=1.000$, $e^{-1}=0.368$ and $\ln 4.086=1.408$.)}
  \item \pp{4} Count the parameters of a fully connected network with $p=4$ inputs, one hidden
    layer with $k=3$ units and an output layer with $K=2$ units, where every hidden and output
    unit has a bias. Show the count separately for the two weight layers.
  \item \pp{2} A convolutional layer receives a $32\times 32$ image and uses filters of size
    $F=5$ with padding $P=2$ and stride $S=1$. Compute the spatial size of the output feature
    map from $\bigl(W-F+2P\bigr)/S+1$.
\end{enumerate}

\ifdefined\withsolutions
\begin{solutionbox}
\textbf{(a) Forward pass.}
Pre-activations:
\[
  z_1=-1+1\cdot 1+2\cdot 2=-1+1+4=4,\qquad
  z_2=-3-2\cdot 1+1\cdot 2=-3-2+2=-3 .
\]
ReLU activations:
\[
  A_1=g(4)=\max\{0,4\}=4,\qquad A_2=g(-3)=\max\{0,-3\}=0 .
\]
Output:
\[
  f(x)=1+2\cdot 4+4\cdot 0=1+8+0=\mathbf{9}.
\]
The ReLU \emph{switched the second unit off}: its negative pre-activation is clipped to $0$,
so unit 2 contributes nothing for this input (but would for other inputs). This thresholding
is the source of non-linearity: without a non-linear $g$, the network would collapse into a
single linear function of $x$, no matter how many layers it has.

\medskip
\textbf{(b) Softmax and cross-entropy.}
Exponentials: $e^{1.0}=2.718$, $e^{0.0}=1.000$, $e^{-1.0}=0.368$; their sum is
$2.718+1.000+0.368=4.086$. Softmax probabilities:
\[
  \hat{p}_1=\frac{2.718}{4.086}=0.665,\qquad
  \hat{p}_2=\frac{1.000}{4.086}=0.245,\qquad
  \hat{p}_3=\frac{0.368}{4.086}=0.090,
\]
which sum to $1.000$. Cross-entropy loss for true class 1:
\[
  -\ln\hat{p}_1=-\ln\frac{e^{1}}{4.086}=-\bigl(1-\ln 4.086\bigr)=1.408-1=\mathbf{0.408}.
\]
(Equivalently $-\ln 0.665\approx 0.408$; a perfect prediction $\hat{p}_1=1$ would give loss 0.)

\medskip
\textbf{(c) Parameter count for $4\to 3\to 2$.}
Input $\to$ hidden: $4\cdot 3=12$ weights $+\,3$ biases $=15$.
Hidden $\to$ output: $3\cdot 2=6$ weights $+\,2$ biases $=8$.
Total: $15+8=\mathbf{23}$ parameters.

\medskip
\textbf{(d) Convolution output size.}
\[
  \frac{W-F+2P}{S}+1=\frac{32-5+2\cdot 2}{1}+1=\frac{31}{1}+1=32,
\]
so the output feature map is $\mathbf{32\times 32}$ --- padding $P=2$ with $F=5$, $S=1$ is
exactly the ``same'' convolution that preserves the spatial size.
\end{solutionbox}
\fi

% ============================================================
\problem{6: Multiple testing}{20}

\begin{enumerate}[label=(\alph*),leftmargin=7mm,itemsep=2.5mm]
  \item \pp{4} A lab performs $m=20$ \emph{independent} hypothesis tests, each at level
    $\alpha=0.05$, and suppose all 20 null hypotheses are true. Define the family-wise error
    rate (FWER) and compute it for this situation
    \emph{(use $(0.95)^{20}=0.3585$)}. Interpret the result in one sentence.
  \item \pp{7} A research group screens $m=6$ candidate biomarkers; the ordered $p$-values are
    \begin{center}
    \begin{tabular}{lcccccc}
    \toprule
    $j$ & 1 & 2 & 3 & 4 & 5 & 6\\
    \midrule
    $p_{(j)}$ & 0.001 & 0.004 & 0.011 & 0.030 & 0.045 & 0.230\\
    \bottomrule
    \end{tabular}
    \end{center}
    Control the FWER at $\alpha=0.05$ (i) with \emph{Bonferroni} and (ii) with the
    \emph{Holm step-down} procedure. Give the relevant thresholds, state for each method
    exactly which hypotheses are rejected, and explain why Holm can never reject fewer
    hypotheses than Bonferroni.
  \item \pp{6} Apply the \emph{Benjamini--Hochberg} (BH) procedure with FDR level $q=0.05$ to
    the same six $p$-values: give the thresholds $jq/m$, determine which hypotheses are
    rejected, and state precisely what ``FDR control at $q=0.05$'' guarantees for the
    resulting set of rejections.
  \item \pp{3} A genomics group screens $20{,}000$ genes in an \emph{exploratory} study; the
    most promising hits will be re-tested in an independent follow-up experiment. Should they
    control the FWER or the FDR? Justify your recommendation.
\end{enumerate}

\ifdefined\withsolutions
\begin{solutionbox}
\textbf{(a) FWER for $m=20$ independent tests.}
$\text{FWER}=\Pr(\text{at least one Type I error})=\Pr(V\ge 1)$, where $V$ is the number of
true nulls that are (falsely) rejected. With independence and all nulls true:
\[
  \text{FWER}=1-\Pr(\text{no false rejection})=1-(1-0.05)^{20}=1-(0.95)^{20}
  =1-0.3585=\mathbf{0.6415}.
\]
Interpretation: even though each single test is performed at the $5\%$ level, there is a
$64\%$ chance of at least one false discovery among the 20 tests --- testing many hypotheses
without adjustment almost guarantees spurious findings.

\medskip
\textbf{(b) Bonferroni and Holm at $\alpha=0.05$.}
\emph{Bonferroni:} reject $H_{(j)}$ iff $p_{(j)}\le\alpha/m=0.05/6=0.00833$.
$p_{(1)}=0.001\le 0.00833$ \checkmark, $p_{(2)}=0.004\le 0.00833$ \checkmark,
$p_{(3)}=0.011>0.00833$ $\times$.
$\Rightarrow$ Bonferroni rejects $H_{(1)}$ and $H_{(2)}$: \textbf{2 rejections}.

\emph{Holm:} compare $p_{(j)}$ with $\alpha/(m-j+1)$ step by step and stop at the first failure:
\begin{center}
\begin{tabular}{lcccccc}
\toprule
$j$ & 1 & 2 & 3 & 4 & 5 & 6\\
\midrule
$p_{(j)}$ & 0.001 & 0.004 & 0.011 & 0.030 & 0.045 & 0.230\\
$\alpha/(m-j+1)$ & 0.00833 & 0.0100 & 0.0125 & 0.0167 & 0.0250 & 0.0500\\
reject? & \checkmark & \checkmark & \checkmark & stop & --- & ---\\
\bottomrule
\end{tabular}
\end{center}
$0.001\le 0.00833$, $0.004\le 0.0100$, $0.011\le 0.0125$, but $0.030>0.0167$: stop.
$\Rightarrow$ Holm rejects $H_{(1)},H_{(2)},H_{(3)}$: \textbf{3 rejections}.

\emph{Why Holm $\supseteq$ Bonferroni:} Holm's thresholds $\alpha/(m-j+1)$ are $\ge\alpha/m$
for every $j$ (with equality only at $j=1$), so every $p$-value that passes Bonferroni also
passes Holm's less stringent later hurdles --- Holm is uniformly more powerful, yet still
controls the FWER at $\alpha$ without any independence assumption.

\medskip
\textbf{(c) Benjamini--Hochberg at $q=0.05$.}
Thresholds $jq/m=j\cdot 0.05/6$:
\begin{center}
\begin{tabular}{lcccccc}
\toprule
$j$ & 1 & 2 & 3 & 4 & 5 & 6\\
\midrule
$p_{(j)}$ & 0.001 & 0.004 & 0.011 & 0.030 & 0.045 & 0.230\\
$jq/m$ & 0.00833 & 0.0167 & 0.0250 & 0.0333 & 0.0417 & 0.0500\\
$p_{(j)}\le jq/m$? & yes & yes & yes & \textbf{yes} & no & no\\
\bottomrule
\end{tabular}
\end{center}
$L=$ largest $j$ with $p_{(j)}\le jq/m$: $j=4$ (since $0.030\le 0.0333$, while
$0.045>0.0417$ and $0.230>0.0500$). BH rejects \emph{all} $H_{(j)}$ with $j\le L$:
$H_{(1)},\dots,H_{(4)}$: \textbf{4 rejections}.
Guarantee: $\text{FDR}=E\bigl[V/\max(R,1)\bigr]\le q=0.05$ --- \emph{on average} at most
$5\%$ of the rejected hypotheses are false discoveries. It does \emph{not} guarantee that at
most $5\%$ of these particular 4 rejections are false, nor does it bound the probability of
making any false rejection (that would be FWER control).

\medskip
\textbf{(d) Exploratory genomics screening.}
Control the \textbf{FDR}. With $m=20{,}000$ tests, FWER control (e.g.\ Bonferroni at
$0.05/20{,}000=2.5\cdot 10^{-6}$) guards against even a single false positive and therefore
has extremely low power --- almost no true signals would survive. In an exploratory screen
the cost of a few false leads is low, because every hit is validated in an independent
follow-up experiment; what matters is that the \emph{list} of discoveries is mostly reliable.
FDR control at, say, $q=0.05$--$0.2$ delivers exactly that: a rich list of candidates of
which on average at most a known fraction is false.
\end{solutionbox}
\fi

% ============================================================
\problem{7: Cumulative essentials (Chapters 2--6)}{20}

\begin{enumerate}[label=(\alph*),leftmargin=7mm,itemsep=2.5mm]
  \item \pp{4} True or false --- justify briefly: ``A more flexible statistical learning method
    always yields a lower \emph{test} error than a less flexible one.''
  \item \pp{4} A logistic regression for loan default,
    $\log\bigl(\frac{p(x)}{1-p(x)}\bigr)=\beta_0+\beta_1 x$ with $x=$ number of missed
    payments, yields $\hat{\beta}_1=0.7$. Interpret the effect of \emph{one additional missed
    payment} on the \emph{odds} of default \emph{(use $e^{0.7}=2.014$)}. Why can the effect on
    the \emph{probability} of default not be stated as a single number?
  \item \pp{4} Give one reason to prefer $k$-fold cross-validation with $k=5$ or $10$ over
    \emph{LOOCV}, and one reason to prefer it over a \emph{single validation split}.
  \item \pp{2} You have $p=2{,}000$ predictors but expect only a handful to be truly related to
    the response. Ridge or lasso --- which penalty is more appropriate, and why?
  \item \pp{6} Match each scenario to exactly one method from
    \{\emph{linear regression}, \emph{lasso}, \emph{logistic regression},
    \emph{random forest}\} (each method used exactly once) and justify each choice in one
    sentence:
    \begin{enumerate}[label=(\roman*),nosep,topsep=1mm]
      \item Estimate how much an additional square metre adds to the price of an apartment,
        with a confidence interval, from $n=2{,}000$ sales.
      \item Predict a continuous drug response from $p=5{,}000$ gene-expression measurements
        for $n=100$ patients; clinicians want a short list of relevant genes.
      \item Predict whether a loan applicant will default (yes/no) and report to the regulator
        how each applicant characteristic changes the odds of default.
      \item Predict machine failures from hundreds of sensor variables with strong non-linear
        effects and interactions; only predictive accuracy matters.
    \end{enumerate}
\end{enumerate}

\ifdefined\withsolutions
\begin{solutionbox}
\textbf{(a) False.}
Expected test error $=$ variance $+$ bias$^2$ $+$ irreducible error. More flexibility reduces
bias but \emph{increases} variance; once the variance increase outweighs the bias reduction,
the test error rises again --- the characteristic U-shape. So a more flexible method wins only
up to the sweet spot; e.g.\ when the true $f$ is nearly linear or $n$ is small, a flexible
method overfits and a simple linear fit has \emph{lower} test error. (Only the \emph{training}
error always decreases with flexibility.)

\medskip
\textbf{(b) Odds ratio.}
One additional missed payment increases the log-odds by $\hat{\beta}_1=0.7$, i.e.\ it
\emph{multiplies the odds} of default by $e^{0.7}=2.014$ --- the odds roughly double,
regardless of the starting value of $x$.
The effect on the \emph{probability} depends on where the observation sits on the S-shaped
logistic curve: the same change in log-odds shifts $p(x)$ a lot near $p\approx 0.5$ but very
little when $p$ is close to $0$ or $1$ --- hence no single number describes it.

\medskip
\textbf{(c) $k$-fold CV with $k=5$ or $10$.}
\emph{vs.\ LOOCV:} the $n$ LOOCV fits are trained on almost identical data sets, so their
errors are highly positively correlated; averaging strongly correlated quantities reduces
variance little, so the LOOCV estimate has \emph{higher variance} than 5- or 10-fold CV
(it is also computationally expensive: $n$ fits instead of $k$, absent shortcut formulas).
\emph{vs.\ a single validation split:} the validation estimate depends strongly on the random
split and overestimates the test error because only part of the data is used for training
(\emph{bias}); $k$-fold CV trains on $(k-1)/k$ of the data and averages over $k$ folds, giving
a less biased and more stable estimate.

\medskip
\textbf{(d) Lasso.}
The $\ell_1$ penalty sets many coefficients \emph{exactly to zero} and thus performs variable
selection --- ideal when the truth is sparse (a handful of signals among 2{,}000 predictors),
and it yields an interpretable small model. Ridge shrinks coefficients but keeps all
$p=2{,}000$ predictors in the model and tends to lose when most true coefficients are zero.

\medskip
\textbf{(e) Matching.}
\begin{itemize}[nosep,leftmargin=5mm]
  \item (i) \textbf{Linear regression}: quantitative response, goal is \emph{inference}
    (effect size with confidence interval) --- exactly what OLS with standard errors provides.
  \item (ii) \textbf{Lasso}: $p\gg n$ makes OLS infeasible; the $\ell_1$ penalty regularises
    and selects a short list of genes.
  \item (iii) \textbf{Logistic regression}: binary response with interpretable coefficients
    ($e^{\hat\beta_j}$ = odds ratios) --- suitable for regulatory reporting.
  \item (iv) \textbf{Random forest}: captures non-linearities and interactions automatically,
    is accurate and robust with many mixed predictors; interpretability is not required.
\end{itemize}
\end{solutionbox}
\fi

\vspace{6mm}
\begin{center}
\rule{0.35\textwidth}{0.4pt}\\[1mm]
\textit{End of the examination --- good luck!}
\end{center}

\end{document}
"
% ============================================================
\documentclass[11pt,a4paper]{article}
\usepackage[utf8]{inputenc}
\usepackage[T1]{fontenc}
\usepackage[english]{babel}
\usepackage[margin=2.2cm]{geometry}
\usepackage{amsmath,amssymb}
\usepackage{booktabs,array}
\usepackage{enumitem}
\usepackage{xcolor}
\usepackage{tcolorbox}
\tcbuselibrary{skins,breakable}

\setlength{\parindent}{0pt}

% Teal solution box (only shown when \withsolutions is defined)
\newtcolorbox{solutionbox}{enhanced,breakable,colback=teal!5,colframe=teal!55!black,
  boxrule=0pt,leftrule=2.5pt,arc=2pt,fonttitle=\bfseries,title=Solution,
  left=2mm,right=2mm,top=1mm,bottom=1.5mm}

\newcommand{\problem}[2]{\par\vspace{7mm}\noindent{\large\textbf{Problem #1}}%
  \hfill{\large\textbf{(#2 points)}}\par\nopagebreak\vspace{0.6mm}%
  \noindent\rule{\textwidth}{0.6pt}\par\nopagebreak\vspace{2.5mm}}
\newcommand{\pp}[1]{\textbf{[#1~P]}\enspace}

\begin{document}

% ------------------------------------------------------------
% Header block
% ------------------------------------------------------------
\begin{center}
  {\Large\textbf{HSBI --- Bielefeld University of Applied Sciences and Arts}}\\[1.2mm]
  {\large Quantitative Research Methods}\\[1mm]
  {\normalsize Summer Semester 2026 \quad$\cdot$\quad Examiner: Prof.\ Dr.\ Christoph Weisser}\\[4.5mm]
  {\LARGE\textbf{Final Mock Examination (Lectures 1--12)}}\\[1.2mm]
  {\normalsize Comprehensive coverage of ISLP Chapters 2--8, 10 and 13,\\
   with emphasis on Chapters 7 (splines/GAMs), 8 (trees), 10 (deep learning) and 13 (multiple testing)}%
  \ifdefined\withsolutions\\[2.2mm]{\large\textcolor{teal!55!black}{\textbf{--- Version with detailed solutions ---}}}\fi
\end{center}

\vspace{2.5mm}
\begin{center}
\begin{tabular}{@{}ll@{}}
  \textbf{Duration:} & \textbf{120 minutes}\\
  \textbf{Total credit:} & \textbf{120 points}\\
  \textbf{Allowed aids:} & non-programmable calculator; one handwritten A4 sheet of notes (both sides)\\
\end{tabular}
\end{center}

\vspace{3mm}
Name: \rule[-1mm]{0.38\textwidth}{0.4pt}\hfill Matriculation no.: \rule[-1mm]{0.24\textwidth}{0.4pt}

\vspace{4.5mm}
\begin{center}
\renewcommand{\arraystretch}{1.35}
\begin{tabular}{|l|c|c|c|c|c|c|c|c|}
\hline
\textbf{Problem} & \textbf{P1} & \textbf{P2} & \textbf{P3} & \textbf{P4} & \textbf{P5} & \textbf{P6} & \textbf{P7} & \textbf{Total}\\\hline
Maximum points & 16 & 8 & 20 & 16 & 20 & 20 & 20 & \textbf{120}\\\hline
Achieved points & \hspace{8mm} & \hspace{8mm} & \hspace{8mm} & \hspace{8mm} & \hspace{8mm} & \hspace{8mm} & \hspace{8mm} & \hspace{10mm}\\\hline
\end{tabular}
\end{center}

\vspace{2.5mm}
\textbf{Instructions}
\begin{itemize}[nosep,leftmargin=5.5mm]
  \item Justify all answers. Numerical results without a derivation or a brief reasoning receive no credit.
  \item Round final numerical results to 3 decimal places unless stated otherwise. Small deviations caused by carrying rounded intermediate results are accepted.
  \item All numerical constants you need (powers of $e$, logarithms, $(0.95)^{20}$, $t$-quantiles) are provided in the problems.
  \item The number of points per problem roughly equals the recommended working time in minutes.
\end{itemize}
\vspace{1mm}\noindent\rule{\textwidth}{0.6pt}

% ============================================================
\problem{1: Moving beyond linearity --- splines}{16}

\begin{enumerate}[label=(\alph*),leftmargin=7mm,itemsep=2.5mm]
  \item \pp{4} A \emph{cubic spline} is a piecewise cubic polynomial that is continuous and has
    continuous first and second derivatives at each knot. How many degrees of freedom (free
    parameters, counting the intercept) does a cubic spline with $K$ interior knots use?
    Justify your answer by counting the parameters of the polynomial pieces and the constraints
    at the knots, and evaluate the result for $K=4$.
  \item \pp{4} How many degrees of freedom does a \emph{natural} cubic spline with the same $K$
    knots use (again evaluate for $K=4$)? Which additional constraint does a natural spline
    impose, in which region of the predictor space does it act, and why is this constraint
    useful in practice?
  \item \pp{4} A \emph{smoothing spline} is the function $g$ that minimises
    \[
      \sum_{i=1}^{n}\bigl(y_i-g(x_i)\bigr)^2+\lambda\int g''(t)^2\,dt .
    \]
    Explain the role of the two terms, and describe the behaviour of the solution $\hat{g}$ in
    the two limiting cases $\lambda\to 0$ and $\lambda\to\infty$.
  \item \pp{4} Write down the \emph{truncated-power basis} representation of a cubic spline with
    a single knot at $\xi$, define the truncated-power basis function explicitly, and explain in
    one sentence why adding this basis function does not destroy the smoothness of the fit at
    $\xi$.
\end{enumerate}

\ifdefined\withsolutions
\begin{solutionbox}
\textbf{(a) Degrees of freedom of a cubic spline.}
$K$ interior knots split the range of $X$ into $K+1$ intervals, each carrying a cubic
polynomial with $4$ coefficients: $4(K+1)=4K+4$ parameters. At each knot, continuity of
$g$, $g'$ and $g''$ imposes $3$ constraints, i.e.\ $3K$ constraints in total. Hence
\[
  \mathrm{df}=4(K+1)-3K=K+4 .
\]
For $K=4$: $\mathrm{df}=4+4=\mathbf{8}$ free parameters (equivalently: intercept, $x$, $x^2$,
$x^3$ plus one truncated-power basis function per knot $=4+K$).

\medskip
\textbf{(b) Natural cubic spline.}
A natural spline additionally requires the function to be \emph{linear beyond the boundary
knots}, i.e.\ in the two outer regions $X<\xi_1$ and $X>\xi_K$. Forcing the quadratic and cubic
terms to vanish there removes $2$ parameters per boundary region, i.e.\ $4$ in total:
\[
  \mathrm{df}=(K+4)-4=K,\qquad K=4:\ \mathrm{df}=\mathbf{4}.
\]
Usefulness: polynomials and unconstrained splines are notoriously erratic near the boundaries,
where data are sparse; the linearity constraint stabilises the fit there
(\emph{lower variance} of predictions at and beyond the boundary), at the price of a small
amount of additional bias.

\medskip
\textbf{(c) Smoothing spline.}
The first term (RSS) rewards fidelity to the data; the second term penalises roughness, since
$g''$ measures the curvature (wiggliness) of $g$, and $\lambda\ge 0$ is the tuning parameter
that trades the two off.
\emph{$\lambda\to 0$:} the penalty disappears; $\hat{g}$ can be made arbitrarily wiggly and
\emph{interpolates} the training observations exactly (zero training RSS, severe overfitting;
effective df $\to n$).
\emph{$\lambda\to\infty$:} any curvature is infinitely expensive, so $g''\equiv 0$ is enforced;
$\hat{g}$ becomes a straight line and equals the \emph{least-squares regression line}
(effective df $\to 2$). Between the extremes, $\lambda$ (chosen e.g.\ by LOOCV) controls the
effective degrees of freedom, which decrease smoothly from $n$ to $2$.

\medskip
\textbf{(d) Truncated-power basis with one knot $\xi$.}
\[
  f(x)=\beta_0+\beta_1x+\beta_2x^2+\beta_3x^3+\beta_4\,(x-\xi)^3_+,
  \qquad
  (x-\xi)^3_+=\begin{cases}(x-\xi)^3, & x>\xi\\[0.5mm] 0, & x\le\xi\end{cases}
\]
The function $(x-\xi)^3_+$ and its first and second derivatives are all $0$ at $x=\xi$, so
adding it changes only the third derivative at the knot; the fit remains continuous with
continuous first and second derivatives --- exactly the defining smoothness of a cubic spline.
\end{solutionbox}
\fi

% ============================================================
\problem{2: Generalised additive models}{8}

\begin{enumerate}[label=(\alph*),leftmargin=7mm,itemsep=2.5mm]
  \item \pp{3} Write down the general form of a \emph{generalised additive model} (GAM) for a
    quantitative response $Y$ and predictors $X_1,\dots,X_p$, and explain in one sentence how it
    extends multiple linear regression.
  \item \pp{3} State one important \emph{advantage} of GAMs that concerns the interpretation of
    the fitted functions $f_j$, and the main structural \emph{limitation} of GAMs --- including
    how this limitation can be repaired.
  \item \pp{2} Describe in one or two sentences how the \emph{backfitting} algorithm fits a GAM.
\end{enumerate}

\ifdefined\withsolutions
\begin{solutionbox}
\textbf{(a) Model form.}
\[
  Y=\beta_0+f_1(X_1)+f_2(X_2)+\dots+f_p(X_p)+\varepsilon .
\]
Multiple linear regression is the special case $f_j(X_j)=\beta_jX_j$; a GAM replaces each linear
component by a smooth non-linear function $f_j$ (e.g.\ a natural or smoothing spline, or local
regression), while keeping the \emph{additive} structure.

\medskip
\textbf{(b) Advantage and limitation.}
\emph{Advantage (interpretability):} because the model is additive, the contribution of each
$X_j$ to $Y$ is captured by its own function $f_j$, which can be examined and plotted
individually while the other variables are held fixed --- non-linear effects remain as easy to
read as in a coefficient table, one panel per predictor.
\emph{Limitation:} additivity excludes \emph{interactions}: the effect of $X_j$ cannot depend
on the value of $X_k$. Repair: add interaction terms manually, e.g.\ a product term
$X_j\times X_k$ or a low-dimensional smooth surface $f_{jk}(X_j,X_k)$ as an extra component.

\medskip
\textbf{(c) Backfitting.}
Cycle through the predictors and repeatedly re-fit each $f_j$ (by a smoother, e.g.\ a spline)
to the \emph{partial residuals} $r_i=y_i-\beta_0-\sum_{k\ne j}f_k(x_{ik})$, holding all other
fitted functions fixed; iterate until the $f_j$ stop changing.
\end{solutionbox}
\fi

% ============================================================
\problem{3: Regression and classification trees by hand}{20}

An online platform records for $n=6$ products the advertising budget $x$ (in 1{,}000 EUR) and
the weekly sales $y$ (in 1{,}000 units):

\begin{center}
\begin{tabular}{lcccccc}
\toprule
Product $i$ & 1 & 2 & 3 & 4 & 5 & 6\\
\midrule
Budget $x_i$ & 1 & 2 & 3 & 4 & 5 & 6\\
Sales $y_i$ & 2 & 4 & 3 & 3 & 8 & 10\\
\bottomrule
\end{tabular}
\end{center}

A regression tree with a \emph{single} split at cutpoint $s$ partitions the data into
$R_1=\{i: x_i<s\}$ and $R_2=\{i: x_i\ge s\}$ and predicts the region mean
$\hat{y}_{R_m}$ in each region.

\begin{enumerate}[label=(\alph*),leftmargin=7mm,itemsep=2.5mm]
  \item \pp{8} For each of the two candidate cutpoints $s=2.5$ and $s=4.5$, determine the two
    regions, the region means and the resulting total residual sum of squares
    $\operatorname{RSS}(s)=\sum_{m=1}^{2}\sum_{i\in R_m}(y_i-\hat{y}_{R_m})^2$.
    Which cutpoint does recursive binary splitting choose, and why?
  \item \pp{2} Using the chosen split, predict the weekly sales for a new product with
    advertising budget $x=5$.
  \item \pp{6} A \emph{classification} tree for a different task has a terminal node containing
    8 training observations: 6 of class ``Yes'' and 2 of class ``No''. Compute the class
    proportions, the Gini index $G=\sum_{k}\hat{p}_{k}(1-\hat{p}_{k})$ and the entropy
    $D=-\sum_{k}\hat{p}_{k}\ln\hat{p}_{k}$ of this node
    \emph{(use $\ln 0.75=-0.2877$ and $\ln 0.25=-1.3863$)}.
    Which class does the node predict?
  \item \pp{4} In \emph{cost-complexity pruning} one minimises, for a subtree
    $T\subseteq T_0$ of the large tree $T_0$,
    \[
      \sum_{m=1}^{|T|}\ \sum_{i:\,x_i\in R_m}\bigl(y_i-\hat{y}_{R_m}\bigr)^2+\alpha\,|T| .
    \]
    Explain the role of the tuning parameter $\alpha$: what happens for $\alpha=0$, what happens
    as $\alpha$ increases, and how is $\alpha$ chosen in practice?
\end{enumerate}

\ifdefined\withsolutions
\begin{solutionbox}
\textbf{(a) Evaluating the two cutpoints.}
Root node for reference: $\bar{y}=30/6=5$.

\emph{Cutpoint $s=2.5$:} $R_1=\{1,2\}$ with $y$-values $\{2,4\}$, mean $\hat{y}_{R_1}=3$;
$R_2=\{3,4,5,6\}$ with $y$-values $\{3,3,8,10\}$, mean $\hat{y}_{R_2}=24/4=6$.
\[
  \operatorname{RSS}_{R_1}=(2-3)^2+(4-3)^2=2,\quad
  \operatorname{RSS}_{R_2}=(3-6)^2+(3-6)^2+(8-6)^2+(10-6)^2=9+9+4+16=38,
\]
so $\operatorname{RSS}(2.5)=2+38=\mathbf{40}$.

\emph{Cutpoint $s=4.5$:} $R_1=\{1,2,3,4\}$ with $y$-values $\{2,4,3,3\}$, mean
$\hat{y}_{R_1}=12/4=3$; $R_2=\{5,6\}$ with $y$-values $\{8,10\}$, mean $\hat{y}_{R_2}=9$.
\[
  \operatorname{RSS}_{R_1}=1+1+0+0=2,\qquad
  \operatorname{RSS}_{R_2}=(8-9)^2+(10-9)^2=2,
\]
so $\operatorname{RSS}(4.5)=2+2=\mathbf{4}$.

Since $4<40$, recursive binary splitting is \emph{greedy} and picks the cutpoint with the
smallest RSS: $\mathbf{s=4.5}$. (It separates the four low-sales products, mean $3$, from the
two high-sales products, mean $9$.)

\medskip
\textbf{(b) Prediction for $x=5$.}
$x=5\ge 4.5$, so the new product falls into $R_2$ and the tree predicts
$\hat{y}=\hat{y}_{R_2}=\mathbf{9}$ (i.e.\ $9{,}000$ units per week).

\medskip
\textbf{(c) Node impurity.}
$\hat{p}_{\text{Yes}}=6/8=0.75$, $\hat{p}_{\text{No}}=2/8=0.25$.
\[
  G=0.75\cdot 0.25+0.25\cdot 0.75=0.1875+0.1875=\mathbf{0.375}.
\]
\[
  D=-\bigl(0.75\ln 0.75+0.25\ln 0.25\bigr)
   =-\bigl(0.75\cdot(-0.2877)+0.25\cdot(-1.3863)\bigr)
   =0.2158+0.3466=\mathbf{0.562}.
\]
The node predicts the majority class, \textbf{Yes} (with estimated probability $0.75$).
Both $G$ and $D$ are small when the node is pure; here the node is moderately impure.

\medskip
\textbf{(d) Role of $\alpha$.}
$\alpha$ is the price per terminal node: it trades off the fit of the subtree (first term)
against its complexity $|T|$.
For $\alpha=0$ the criterion equals the training RSS and the full tree $T_0$ is optimal.
As $\alpha$ increases, additional leaves must ``pay for themselves''; branches whose RSS
reduction is smaller than $\alpha$ are pruned, so the optimal subtree becomes smaller ---
one obtains a \emph{nested sequence} of subtrees as $\alpha$ grows (weakest-link pruning).
In practice $\alpha$ is selected by \emph{cross-validation} (choose the $\alpha$ with the
smallest CV error, then refit on all data), which balances bias against variance instead of
optimising training fit.
\end{solutionbox}
\fi

% ============================================================
\problem{4: Ensembles --- bagging, random forests, boosting}{16}

\begin{enumerate}[label=(\alph*),leftmargin=7mm,itemsep=2.5mm]
  \item \pp{6} Let $\hat{f}_1(x),\dots,\hat{f}_B(x)$ be $B$ identically distributed (but not
    independent) tree predictions at a fixed point $x$, each with variance $\sigma^2$ and with
    pairwise correlation $\rho$. The variance of the bagged average is
    \[
      \operatorname{Var}\Bigl(\tfrac{1}{B}\sum_{b=1}^{B}\hat{f}_b(x)\Bigr)
      =\rho\,\sigma^{2}+\frac{1-\rho}{B}\,\sigma^{2}.
    \]
    (i) Explain briefly why this formula shows that bagging reduces variance, and which term
    limits the reduction. (ii) Evaluate the variance for $B=100$, $\rho=0.6$ and $\sigma^2=4$,
    and compare it with the variance of a single tree. (iii) What value does the variance
    approach as $B\to\infty$?
  \item \pp{4} How do \emph{random forests} modify bagging at each split, and what is the
    recommended subset size $m$ for classification as a function of $p$ (evaluate for $p=25$)?
    Explain, with reference to the formula in (a), why this modification improves on bagging.
  \item \pp{4} Name the three tuning parameters of \emph{boosting} for regression trees and
    state in one sentence each what they control. What is the ``slow learning'' idea behind
    boosting?
  \item \pp{2} What does the \emph{out-of-bag} (OOB) error of a bagged model estimate, and why
    is it available essentially for free?
\end{enumerate}

\ifdefined\withsolutions
\begin{solutionbox}
\textbf{(a) Variance of the bagged average.}
(i) The second term $\frac{1-\rho}{B}\sigma^2$ is driven to $0$ by averaging over many trees
($B$ large), so averaging removes the ``independent part'' of the trees' variability ---
this is the variance reduction of bagging. The first term $\rho\sigma^2$ does \emph{not}
depend on $B$: because the bootstrap trees are grown on overlapping samples of the same data
(and share the same strong predictors), they are positively correlated, and this correlated
part of the variance cannot be averaged away. It limits the reduction.
(ii) For $B=100$, $\rho=0.6$, $\sigma^2=4$:
\[
  \operatorname{Var}=0.6\cdot 4+\frac{1-0.6}{100}\cdot 4
  =2.4+\frac{0.4\cdot 4}{100}=2.4+0.016=\mathbf{2.416},
\]
compared with $\sigma^2=4$ for a single tree --- a reduction of about $40\%$.
(iii) As $B\to\infty$ the variance approaches the floor $\rho\sigma^2=0.6\cdot 4=\mathbf{2.4}$.
(Increasing $B$ further does not overfit; it only stabilises the average.)

\medskip
\textbf{(b) Random forests.}
At each split, only a \emph{random subset} of $m$ of the $p$ predictors is considered as
splitting candidates; a fresh subset is drawn at every split. Recommended for classification:
$m\approx\sqrt{p}$, so for $p=25$: $m=\sqrt{25}=\mathbf{5}$.
Why it helps: with bagging, a dominant predictor is used in the top split of almost every tree,
so the trees look alike and $\rho$ is high. Excluding it from $\approx (p-m)/p$ of the split
searches \emph{decorrelates} the trees, i.e.\ lowers $\rho$ --- and by (a) this lowers the
irreducible floor $\rho\sigma^2$ of the ensemble variance, which mere averaging cannot touch.

\medskip
\textbf{(c) Boosting.}
\begin{itemize}[nosep,leftmargin=5mm]
  \item \emph{Number of trees $B$:} how many small trees are added sequentially; unlike
    bagging, boosting \emph{can} overfit if $B$ is too large, so $B$ is chosen by
    cross-validation.
  \item \emph{Shrinkage $\lambda$} (learning rate, e.g.\ $0.01$ or $0.001$): scales the
    contribution of each new tree and thus the speed of learning; smaller $\lambda$ usually
    needs larger $B$.
  \item \emph{Number of splits $d$ per tree} (interaction depth): the complexity of each tree
    and the interaction order of the model; often even $d=1$ (stumps, additive model) works
    well.
\end{itemize}
\emph{Slow learning:} each tree is fitted to the current \emph{residuals}, and only a small
fraction $\lambda$ of it is added to the model; the model thus improves slowly and exactly in
the regions where it still performs poorly, instead of fitting the data hard in one step ---
which tends to generalise better.

\medskip
\textbf{(d) OOB error.}
Each bootstrap sample leaves out on average about $1/3$ of the observations
(``out-of-bag''). Predicting each observation only with the trees that did \emph{not} use it,
and aggregating these predictions, yields the OOB error --- a valid estimate of the
\emph{test error} of the bagged model. It is free because it reuses the bootstrap side
product; no separate validation set or cross-validation loop is required.
\end{solutionbox}
\fi

% ============================================================
\problem{5: Neural networks by hand}{20}

Consider a tiny feed-forward network for a quantitative response with input
$x=(x_1,x_2)=(1,2)$, one hidden layer with two ReLU units
($g(z)=\max\{0,z\}$) and a linear output unit:
\[
  A_1=g\bigl(-1+1\,x_1+2\,x_2\bigr),\qquad
  A_2=g\bigl(-3-2\,x_1+1\,x_2\bigr),\qquad
  f(x)=1+2\,A_1+4\,A_2 .
\]

\begin{enumerate}[label=(\alph*),leftmargin=7mm,itemsep=2.5mm]
  \item \pp{8} Perform the forward pass step by step: compute both pre-activations, both
    activations $A_1,A_2$ and the network output $f(x)$. Comment briefly on what the ReLU did
    to the second unit and why such non-linear activations are essential.
  \item \pp{6} A classification network with $K=3$ classes produces the logits (output-layer
    pre-activations) $z=(z_1,z_2,z_3)=(1.0,\ 0.0,\ -1.0)$ for one observation whose true class
    is class~1. Compute all three softmax probabilities and the cross-entropy loss
    $-\ln \hat{p}_{\text{true}}$.
    \emph{(Use $e^{1}=2.718$, $e^{0}=1.000$, $e^{-1}=0.368$ and $\ln 4.086=1.408$.)}
  \item \pp{4} Count the parameters of a fully connected network with $p=4$ inputs, one hidden
    layer with $k=3$ units and an output layer with $K=2$ units, where every hidden and output
    unit has a bias. Show the count separately for the two weight layers.
  \item \pp{2} A convolutional layer receives a $32\times 32$ image and uses filters of size
    $F=5$ with padding $P=2$ and stride $S=1$. Compute the spatial size of the output feature
    map from $\bigl(W-F+2P\bigr)/S+1$.
\end{enumerate}

\ifdefined\withsolutions
\begin{solutionbox}
\textbf{(a) Forward pass.}
Pre-activations:
\[
  z_1=-1+1\cdot 1+2\cdot 2=-1+1+4=4,\qquad
  z_2=-3-2\cdot 1+1\cdot 2=-3-2+2=-3 .
\]
ReLU activations:
\[
  A_1=g(4)=\max\{0,4\}=4,\qquad A_2=g(-3)=\max\{0,-3\}=0 .
\]
Output:
\[
  f(x)=1+2\cdot 4+4\cdot 0=1+8+0=\mathbf{9}.
\]
The ReLU \emph{switched the second unit off}: its negative pre-activation is clipped to $0$,
so unit 2 contributes nothing for this input (but would for other inputs). This thresholding
is the source of non-linearity: without a non-linear $g$, the network would collapse into a
single linear function of $x$, no matter how many layers it has.

\medskip
\textbf{(b) Softmax and cross-entropy.}
Exponentials: $e^{1.0}=2.718$, $e^{0.0}=1.000$, $e^{-1.0}=0.368$; their sum is
$2.718+1.000+0.368=4.086$. Softmax probabilities:
\[
  \hat{p}_1=\frac{2.718}{4.086}=0.665,\qquad
  \hat{p}_2=\frac{1.000}{4.086}=0.245,\qquad
  \hat{p}_3=\frac{0.368}{4.086}=0.090,
\]
which sum to $1.000$. Cross-entropy loss for true class 1:
\[
  -\ln\hat{p}_1=-\ln\frac{e^{1}}{4.086}=-\bigl(1-\ln 4.086\bigr)=1.408-1=\mathbf{0.408}.
\]
(Equivalently $-\ln 0.665\approx 0.408$; a perfect prediction $\hat{p}_1=1$ would give loss 0.)

\medskip
\textbf{(c) Parameter count for $4\to 3\to 2$.}
Input $\to$ hidden: $4\cdot 3=12$ weights $+\,3$ biases $=15$.
Hidden $\to$ output: $3\cdot 2=6$ weights $+\,2$ biases $=8$.
Total: $15+8=\mathbf{23}$ parameters.

\medskip
\textbf{(d) Convolution output size.}
\[
  \frac{W-F+2P}{S}+1=\frac{32-5+2\cdot 2}{1}+1=\frac{31}{1}+1=32,
\]
so the output feature map is $\mathbf{32\times 32}$ --- padding $P=2$ with $F=5$, $S=1$ is
exactly the ``same'' convolution that preserves the spatial size.
\end{solutionbox}
\fi

% ============================================================
\problem{6: Multiple testing}{20}

\begin{enumerate}[label=(\alph*),leftmargin=7mm,itemsep=2.5mm]
  \item \pp{4} A lab performs $m=20$ \emph{independent} hypothesis tests, each at level
    $\alpha=0.05$, and suppose all 20 null hypotheses are true. Define the family-wise error
    rate (FWER) and compute it for this situation
    \emph{(use $(0.95)^{20}=0.3585$)}. Interpret the result in one sentence.
  \item \pp{7} A research group screens $m=6$ candidate biomarkers; the ordered $p$-values are
    \begin{center}
    \begin{tabular}{lcccccc}
    \toprule
    $j$ & 1 & 2 & 3 & 4 & 5 & 6\\
    \midrule
    $p_{(j)}$ & 0.001 & 0.004 & 0.011 & 0.030 & 0.045 & 0.230\\
    \bottomrule
    \end{tabular}
    \end{center}
    Control the FWER at $\alpha=0.05$ (i) with \emph{Bonferroni} and (ii) with the
    \emph{Holm step-down} procedure. Give the relevant thresholds, state for each method
    exactly which hypotheses are rejected, and explain why Holm can never reject fewer
    hypotheses than Bonferroni.
  \item \pp{6} Apply the \emph{Benjamini--Hochberg} (BH) procedure with FDR level $q=0.05$ to
    the same six $p$-values: give the thresholds $jq/m$, determine which hypotheses are
    rejected, and state precisely what ``FDR control at $q=0.05$'' guarantees for the
    resulting set of rejections.
  \item \pp{3} A genomics group screens $20{,}000$ genes in an \emph{exploratory} study; the
    most promising hits will be re-tested in an independent follow-up experiment. Should they
    control the FWER or the FDR? Justify your recommendation.
\end{enumerate}

\ifdefined\withsolutions
\begin{solutionbox}
\textbf{(a) FWER for $m=20$ independent tests.}
$\text{FWER}=\Pr(\text{at least one Type I error})=\Pr(V\ge 1)$, where $V$ is the number of
true nulls that are (falsely) rejected. With independence and all nulls true:
\[
  \text{FWER}=1-\Pr(\text{no false rejection})=1-(1-0.05)^{20}=1-(0.95)^{20}
  =1-0.3585=\mathbf{0.6415}.
\]
Interpretation: even though each single test is performed at the $5\%$ level, there is a
$64\%$ chance of at least one false discovery among the 20 tests --- testing many hypotheses
without adjustment almost guarantees spurious findings.

\medskip
\textbf{(b) Bonferroni and Holm at $\alpha=0.05$.}
\emph{Bonferroni:} reject $H_{(j)}$ iff $p_{(j)}\le\alpha/m=0.05/6=0.00833$.
$p_{(1)}=0.001\le 0.00833$ \checkmark, $p_{(2)}=0.004\le 0.00833$ \checkmark,
$p_{(3)}=0.011>0.00833$ $\times$.
$\Rightarrow$ Bonferroni rejects $H_{(1)}$ and $H_{(2)}$: \textbf{2 rejections}.

\emph{Holm:} compare $p_{(j)}$ with $\alpha/(m-j+1)$ step by step and stop at the first failure:
\begin{center}
\begin{tabular}{lcccccc}
\toprule
$j$ & 1 & 2 & 3 & 4 & 5 & 6\\
\midrule
$p_{(j)}$ & 0.001 & 0.004 & 0.011 & 0.030 & 0.045 & 0.230\\
$\alpha/(m-j+1)$ & 0.00833 & 0.0100 & 0.0125 & 0.0167 & 0.0250 & 0.0500\\
reject? & \checkmark & \checkmark & \checkmark & stop & --- & ---\\
\bottomrule
\end{tabular}
\end{center}
$0.001\le 0.00833$, $0.004\le 0.0100$, $0.011\le 0.0125$, but $0.030>0.0167$: stop.
$\Rightarrow$ Holm rejects $H_{(1)},H_{(2)},H_{(3)}$: \textbf{3 rejections}.

\emph{Why Holm $\supseteq$ Bonferroni:} Holm's thresholds $\alpha/(m-j+1)$ are $\ge\alpha/m$
for every $j$ (with equality only at $j=1$), so every $p$-value that passes Bonferroni also
passes Holm's less stringent later hurdles --- Holm is uniformly more powerful, yet still
controls the FWER at $\alpha$ without any independence assumption.

\medskip
\textbf{(c) Benjamini--Hochberg at $q=0.05$.}
Thresholds $jq/m=j\cdot 0.05/6$:
\begin{center}
\begin{tabular}{lcccccc}
\toprule
$j$ & 1 & 2 & 3 & 4 & 5 & 6\\
\midrule
$p_{(j)}$ & 0.001 & 0.004 & 0.011 & 0.030 & 0.045 & 0.230\\
$jq/m$ & 0.00833 & 0.0167 & 0.0250 & 0.0333 & 0.0417 & 0.0500\\
$p_{(j)}\le jq/m$? & yes & yes & yes & \textbf{yes} & no & no\\
\bottomrule
\end{tabular}
\end{center}
$L=$ largest $j$ with $p_{(j)}\le jq/m$: $j=4$ (since $0.030\le 0.0333$, while
$0.045>0.0417$ and $0.230>0.0500$). BH rejects \emph{all} $H_{(j)}$ with $j\le L$:
$H_{(1)},\dots,H_{(4)}$: \textbf{4 rejections}.
Guarantee: $\text{FDR}=E\bigl[V/\max(R,1)\bigr]\le q=0.05$ --- \emph{on average} at most
$5\%$ of the rejected hypotheses are false discoveries. It does \emph{not} guarantee that at
most $5\%$ of these particular 4 rejections are false, nor does it bound the probability of
making any false rejection (that would be FWER control).

\medskip
\textbf{(d) Exploratory genomics screening.}
Control the \textbf{FDR}. With $m=20{,}000$ tests, FWER control (e.g.\ Bonferroni at
$0.05/20{,}000=2.5\cdot 10^{-6}$) guards against even a single false positive and therefore
has extremely low power --- almost no true signals would survive. In an exploratory screen
the cost of a few false leads is low, because every hit is validated in an independent
follow-up experiment; what matters is that the \emph{list} of discoveries is mostly reliable.
FDR control at, say, $q=0.05$--$0.2$ delivers exactly that: a rich list of candidates of
which on average at most a known fraction is false.
\end{solutionbox}
\fi

% ============================================================
\problem{7: Cumulative essentials (Chapters 2--6)}{20}

\begin{enumerate}[label=(\alph*),leftmargin=7mm,itemsep=2.5mm]
  \item \pp{4} True or false --- justify briefly: ``A more flexible statistical learning method
    always yields a lower \emph{test} error than a less flexible one.''
  \item \pp{4} A logistic regression for loan default,
    $\log\bigl(\frac{p(x)}{1-p(x)}\bigr)=\beta_0+\beta_1 x$ with $x=$ number of missed
    payments, yields $\hat{\beta}_1=0.7$. Interpret the effect of \emph{one additional missed
    payment} on the \emph{odds} of default \emph{(use $e^{0.7}=2.014$)}. Why can the effect on
    the \emph{probability} of default not be stated as a single number?
  \item \pp{4} Give one reason to prefer $k$-fold cross-validation with $k=5$ or $10$ over
    \emph{LOOCV}, and one reason to prefer it over a \emph{single validation split}.
  \item \pp{2} You have $p=2{,}000$ predictors but expect only a handful to be truly related to
    the response. Ridge or lasso --- which penalty is more appropriate, and why?
  \item \pp{6} Match each scenario to exactly one method from
    \{\emph{linear regression}, \emph{lasso}, \emph{logistic regression},
    \emph{random forest}\} (each method used exactly once) and justify each choice in one
    sentence:
    \begin{enumerate}[label=(\roman*),nosep,topsep=1mm]
      \item Estimate how much an additional square metre adds to the price of an apartment,
        with a confidence interval, from $n=2{,}000$ sales.
      \item Predict a continuous drug response from $p=5{,}000$ gene-expression measurements
        for $n=100$ patients; clinicians want a short list of relevant genes.
      \item Predict whether a loan applicant will default (yes/no) and report to the regulator
        how each applicant characteristic changes the odds of default.
      \item Predict machine failures from hundreds of sensor variables with strong non-linear
        effects and interactions; only predictive accuracy matters.
    \end{enumerate}
\end{enumerate}

\ifdefined\withsolutions
\begin{solutionbox}
\textbf{(a) False.}
Expected test error $=$ variance $+$ bias$^2$ $+$ irreducible error. More flexibility reduces
bias but \emph{increases} variance; once the variance increase outweighs the bias reduction,
the test error rises again --- the characteristic U-shape. So a more flexible method wins only
up to the sweet spot; e.g.\ when the true $f$ is nearly linear or $n$ is small, a flexible
method overfits and a simple linear fit has \emph{lower} test error. (Only the \emph{training}
error always decreases with flexibility.)

\medskip
\textbf{(b) Odds ratio.}
One additional missed payment increases the log-odds by $\hat{\beta}_1=0.7$, i.e.\ it
\emph{multiplies the odds} of default by $e^{0.7}=2.014$ --- the odds roughly double,
regardless of the starting value of $x$.
The effect on the \emph{probability} depends on where the observation sits on the S-shaped
logistic curve: the same change in log-odds shifts $p(x)$ a lot near $p\approx 0.5$ but very
little when $p$ is close to $0$ or $1$ --- hence no single number describes it.

\medskip
\textbf{(c) $k$-fold CV with $k=5$ or $10$.}
\emph{vs.\ LOOCV:} the $n$ LOOCV fits are trained on almost identical data sets, so their
errors are highly positively correlated; averaging strongly correlated quantities reduces
variance little, so the LOOCV estimate has \emph{higher variance} than 5- or 10-fold CV
(it is also computationally expensive: $n$ fits instead of $k$, absent shortcut formulas).
\emph{vs.\ a single validation split:} the validation estimate depends strongly on the random
split and overestimates the test error because only part of the data is used for training
(\emph{bias}); $k$-fold CV trains on $(k-1)/k$ of the data and averages over $k$ folds, giving
a less biased and more stable estimate.

\medskip
\textbf{(d) Lasso.}
The $\ell_1$ penalty sets many coefficients \emph{exactly to zero} and thus performs variable
selection --- ideal when the truth is sparse (a handful of signals among 2{,}000 predictors),
and it yields an interpretable small model. Ridge shrinks coefficients but keeps all
$p=2{,}000$ predictors in the model and tends to lose when most true coefficients are zero.

\medskip
\textbf{(e) Matching.}
\begin{itemize}[nosep,leftmargin=5mm]
  \item (i) \textbf{Linear regression}: quantitative response, goal is \emph{inference}
    (effect size with confidence interval) --- exactly what OLS with standard errors provides.
  \item (ii) \textbf{Lasso}: $p\gg n$ makes OLS infeasible; the $\ell_1$ penalty regularises
    and selects a short list of genes.
  \item (iii) \textbf{Logistic regression}: binary response with interpretable coefficients
    ($e^{\hat\beta_j}$ = odds ratios) --- suitable for regulatory reporting.
  \item (iv) \textbf{Random forest}: captures non-linearities and interactions automatically,
    is accurate and robust with many mixed predictors; interpretability is not required.
\end{itemize}
\end{solutionbox}
\fi

\vspace{6mm}
\begin{center}
\rule{0.35\textwidth}{0.4pt}\\[1mm]
\textit{End of the examination --- good luck!}
\end{center}

\end{document}
"
% ============================================================
\documentclass[11pt,a4paper]{article}
\usepackage[utf8]{inputenc}
\usepackage[T1]{fontenc}
\usepackage[english]{babel}
\usepackage[margin=2.2cm]{geometry}
\usepackage{amsmath,amssymb}
\usepackage{booktabs,array}
\usepackage{enumitem}
\usepackage{xcolor}
\usepackage{tcolorbox}
\tcbuselibrary{skins,breakable}

\setlength{\parindent}{0pt}

% Teal solution box (only shown when \withsolutions is defined)
\newtcolorbox{solutionbox}{enhanced,breakable,colback=teal!5,colframe=teal!55!black,
  boxrule=0pt,leftrule=2.5pt,arc=2pt,fonttitle=\bfseries,title=Solution,
  left=2mm,right=2mm,top=1mm,bottom=1.5mm}

\newcommand{\problem}[2]{\par\vspace{7mm}\noindent{\large\textbf{Problem #1}}%
  \hfill{\large\textbf{(#2 points)}}\par\nopagebreak\vspace{0.6mm}%
  \noindent\rule{\textwidth}{0.6pt}\par\nopagebreak\vspace{2.5mm}}
\newcommand{\pp}[1]{\textbf{[#1~P]}\enspace}

\begin{document}

% ------------------------------------------------------------
% Header block
% ------------------------------------------------------------
\begin{center}
  {\Large\textbf{HSBI --- Bielefeld University of Applied Sciences and Arts}}\\[1.2mm]
  {\large Quantitative Research Methods}\\[1mm]
  {\normalsize Summer Semester 2026 \quad$\cdot$\quad Examiner: Prof.\ Dr.\ Christoph Weisser}\\[4.5mm]
  {\LARGE\textbf{Final Mock Examination (Lectures 1--12)}}\\[1.2mm]
  {\normalsize Comprehensive coverage of ISLP Chapters 2--8, 10 and 13,\\
   with emphasis on Chapters 7 (splines/GAMs), 8 (trees), 10 (deep learning) and 13 (multiple testing)}%
  \ifdefined\withsolutions\\[2.2mm]{\large\textcolor{teal!55!black}{\textbf{--- Version with detailed solutions ---}}}\fi
\end{center}

\vspace{2.5mm}
\begin{center}
\begin{tabular}{@{}ll@{}}
  \textbf{Duration:} & \textbf{120 minutes}\\
  \textbf{Total credit:} & \textbf{120 points}\\
  \textbf{Allowed aids:} & non-programmable calculator; one handwritten A4 sheet of notes (both sides)\\
\end{tabular}
\end{center}

\vspace{3mm}
Name: \rule[-1mm]{0.38\textwidth}{0.4pt}\hfill Matriculation no.: \rule[-1mm]{0.24\textwidth}{0.4pt}

\vspace{4.5mm}
\begin{center}
\renewcommand{\arraystretch}{1.35}
\begin{tabular}{|l|c|c|c|c|c|c|c|c|}
\hline
\textbf{Problem} & \textbf{P1} & \textbf{P2} & \textbf{P3} & \textbf{P4} & \textbf{P5} & \textbf{P6} & \textbf{P7} & \textbf{Total}\\\hline
Maximum points & 16 & 8 & 20 & 16 & 20 & 20 & 20 & \textbf{120}\\\hline
Achieved points & \hspace{8mm} & \hspace{8mm} & \hspace{8mm} & \hspace{8mm} & \hspace{8mm} & \hspace{8mm} & \hspace{8mm} & \hspace{10mm}\\\hline
\end{tabular}
\end{center}

\vspace{2.5mm}
\textbf{Instructions}
\begin{itemize}[nosep,leftmargin=5.5mm]
  \item Justify all answers. Numerical results without a derivation or a brief reasoning receive no credit.
  \item Round final numerical results to 3 decimal places unless stated otherwise. Small deviations caused by carrying rounded intermediate results are accepted.
  \item All numerical constants you need (powers of $e$, logarithms, $(0.95)^{20}$, $t$-quantiles) are provided in the problems.
  \item The number of points per problem roughly equals the recommended working time in minutes.
\end{itemize}
\vspace{1mm}\noindent\rule{\textwidth}{0.6pt}

% ============================================================
\problem{1: Moving beyond linearity --- splines}{16}

\begin{enumerate}[label=(\alph*),leftmargin=7mm,itemsep=2.5mm]
  \item \pp{4} A \emph{cubic spline} is a piecewise cubic polynomial that is continuous and has
    continuous first and second derivatives at each knot. How many degrees of freedom (free
    parameters, counting the intercept) does a cubic spline with $K$ interior knots use?
    Justify your answer by counting the parameters of the polynomial pieces and the constraints
    at the knots, and evaluate the result for $K=4$.
  \item \pp{4} How many degrees of freedom does a \emph{natural} cubic spline with the same $K$
    knots use (again evaluate for $K=4$)? Which additional constraint does a natural spline
    impose, in which region of the predictor space does it act, and why is this constraint
    useful in practice?
  \item \pp{4} A \emph{smoothing spline} is the function $g$ that minimises
    \[
      \sum_{i=1}^{n}\bigl(y_i-g(x_i)\bigr)^2+\lambda\int g''(t)^2\,dt .
    \]
    Explain the role of the two terms, and describe the behaviour of the solution $\hat{g}$ in
    the two limiting cases $\lambda\to 0$ and $\lambda\to\infty$.
  \item \pp{4} Write down the \emph{truncated-power basis} representation of a cubic spline with
    a single knot at $\xi$, define the truncated-power basis function explicitly, and explain in
    one sentence why adding this basis function does not destroy the smoothness of the fit at
    $\xi$.
\end{enumerate}

\ifdefined\withsolutions
\begin{solutionbox}
\textbf{(a) Degrees of freedom of a cubic spline.}
$K$ interior knots split the range of $X$ into $K+1$ intervals, each carrying a cubic
polynomial with $4$ coefficients: $4(K+1)=4K+4$ parameters. At each knot, continuity of
$g$, $g'$ and $g''$ imposes $3$ constraints, i.e.\ $3K$ constraints in total. Hence
\[
  \mathrm{df}=4(K+1)-3K=K+4 .
\]
For $K=4$: $\mathrm{df}=4+4=\mathbf{8}$ free parameters (equivalently: intercept, $x$, $x^2$,
$x^3$ plus one truncated-power basis function per knot $=4+K$).

\medskip
\textbf{(b) Natural cubic spline.}
A natural spline additionally requires the function to be \emph{linear beyond the boundary
knots}, i.e.\ in the two outer regions $X<\xi_1$ and $X>\xi_K$. Forcing the quadratic and cubic
terms to vanish there removes $2$ parameters per boundary region, i.e.\ $4$ in total:
\[
  \mathrm{df}=(K+4)-4=K,\qquad K=4:\ \mathrm{df}=\mathbf{4}.
\]
Usefulness: polynomials and unconstrained splines are notoriously erratic near the boundaries,
where data are sparse; the linearity constraint stabilises the fit there
(\emph{lower variance} of predictions at and beyond the boundary), at the price of a small
amount of additional bias.

\medskip
\textbf{(c) Smoothing spline.}
The first term (RSS) rewards fidelity to the data; the second term penalises roughness, since
$g''$ measures the curvature (wiggliness) of $g$, and $\lambda\ge 0$ is the tuning parameter
that trades the two off.
\emph{$\lambda\to 0$:} the penalty disappears; $\hat{g}$ can be made arbitrarily wiggly and
\emph{interpolates} the training observations exactly (zero training RSS, severe overfitting;
effective df $\to n$).
\emph{$\lambda\to\infty$:} any curvature is infinitely expensive, so $g''\equiv 0$ is enforced;
$\hat{g}$ becomes a straight line and equals the \emph{least-squares regression line}
(effective df $\to 2$). Between the extremes, $\lambda$ (chosen e.g.\ by LOOCV) controls the
effective degrees of freedom, which decrease smoothly from $n$ to $2$.

\medskip
\textbf{(d) Truncated-power basis with one knot $\xi$.}
\[
  f(x)=\beta_0+\beta_1x+\beta_2x^2+\beta_3x^3+\beta_4\,(x-\xi)^3_+,
  \qquad
  (x-\xi)^3_+=\begin{cases}(x-\xi)^3, & x>\xi\\[0.5mm] 0, & x\le\xi\end{cases}
\]
The function $(x-\xi)^3_+$ and its first and second derivatives are all $0$ at $x=\xi$, so
adding it changes only the third derivative at the knot; the fit remains continuous with
continuous first and second derivatives --- exactly the defining smoothness of a cubic spline.
\end{solutionbox}
\fi

% ============================================================
\problem{2: Generalised additive models}{8}

\begin{enumerate}[label=(\alph*),leftmargin=7mm,itemsep=2.5mm]
  \item \pp{3} Write down the general form of a \emph{generalised additive model} (GAM) for a
    quantitative response $Y$ and predictors $X_1,\dots,X_p$, and explain in one sentence how it
    extends multiple linear regression.
  \item \pp{3} State one important \emph{advantage} of GAMs that concerns the interpretation of
    the fitted functions $f_j$, and the main structural \emph{limitation} of GAMs --- including
    how this limitation can be repaired.
  \item \pp{2} Describe in one or two sentences how the \emph{backfitting} algorithm fits a GAM.
\end{enumerate}

\ifdefined\withsolutions
\begin{solutionbox}
\textbf{(a) Model form.}
\[
  Y=\beta_0+f_1(X_1)+f_2(X_2)+\dots+f_p(X_p)+\varepsilon .
\]
Multiple linear regression is the special case $f_j(X_j)=\beta_jX_j$; a GAM replaces each linear
component by a smooth non-linear function $f_j$ (e.g.\ a natural or smoothing spline, or local
regression), while keeping the \emph{additive} structure.

\medskip
\textbf{(b) Advantage and limitation.}
\emph{Advantage (interpretability):} because the model is additive, the contribution of each
$X_j$ to $Y$ is captured by its own function $f_j$, which can be examined and plotted
individually while the other variables are held fixed --- non-linear effects remain as easy to
read as in a coefficient table, one panel per predictor.
\emph{Limitation:} additivity excludes \emph{interactions}: the effect of $X_j$ cannot depend
on the value of $X_k$. Repair: add interaction terms manually, e.g.\ a product term
$X_j\times X_k$ or a low-dimensional smooth surface $f_{jk}(X_j,X_k)$ as an extra component.

\medskip
\textbf{(c) Backfitting.}
Cycle through the predictors and repeatedly re-fit each $f_j$ (by a smoother, e.g.\ a spline)
to the \emph{partial residuals} $r_i=y_i-\beta_0-\sum_{k\ne j}f_k(x_{ik})$, holding all other
fitted functions fixed; iterate until the $f_j$ stop changing.
\end{solutionbox}
\fi

% ============================================================
\problem{3: Regression and classification trees by hand}{20}

An online platform records for $n=6$ products the advertising budget $x$ (in 1{,}000 EUR) and
the weekly sales $y$ (in 1{,}000 units):

\begin{center}
\begin{tabular}{lcccccc}
\toprule
Product $i$ & 1 & 2 & 3 & 4 & 5 & 6\\
\midrule
Budget $x_i$ & 1 & 2 & 3 & 4 & 5 & 6\\
Sales $y_i$ & 2 & 4 & 3 & 3 & 8 & 10\\
\bottomrule
\end{tabular}
\end{center}

A regression tree with a \emph{single} split at cutpoint $s$ partitions the data into
$R_1=\{i: x_i<s\}$ and $R_2=\{i: x_i\ge s\}$ and predicts the region mean
$\hat{y}_{R_m}$ in each region.

\begin{enumerate}[label=(\alph*),leftmargin=7mm,itemsep=2.5mm]
  \item \pp{8} For each of the two candidate cutpoints $s=2.5$ and $s=4.5$, determine the two
    regions, the region means and the resulting total residual sum of squares
    $\operatorname{RSS}(s)=\sum_{m=1}^{2}\sum_{i\in R_m}(y_i-\hat{y}_{R_m})^2$.
    Which cutpoint does recursive binary splitting choose, and why?
  \item \pp{2} Using the chosen split, predict the weekly sales for a new product with
    advertising budget $x=5$.
  \item \pp{6} A \emph{classification} tree for a different task has a terminal node containing
    8 training observations: 6 of class ``Yes'' and 2 of class ``No''. Compute the class
    proportions, the Gini index $G=\sum_{k}\hat{p}_{k}(1-\hat{p}_{k})$ and the entropy
    $D=-\sum_{k}\hat{p}_{k}\ln\hat{p}_{k}$ of this node
    \emph{(use $\ln 0.75=-0.2877$ and $\ln 0.25=-1.3863$)}.
    Which class does the node predict?
  \item \pp{4} In \emph{cost-complexity pruning} one minimises, for a subtree
    $T\subseteq T_0$ of the large tree $T_0$,
    \[
      \sum_{m=1}^{|T|}\ \sum_{i:\,x_i\in R_m}\bigl(y_i-\hat{y}_{R_m}\bigr)^2+\alpha\,|T| .
    \]
    Explain the role of the tuning parameter $\alpha$: what happens for $\alpha=0$, what happens
    as $\alpha$ increases, and how is $\alpha$ chosen in practice?
\end{enumerate}

\ifdefined\withsolutions
\begin{solutionbox}
\textbf{(a) Evaluating the two cutpoints.}
Root node for reference: $\bar{y}=30/6=5$.

\emph{Cutpoint $s=2.5$:} $R_1=\{1,2\}$ with $y$-values $\{2,4\}$, mean $\hat{y}_{R_1}=3$;
$R_2=\{3,4,5,6\}$ with $y$-values $\{3,3,8,10\}$, mean $\hat{y}_{R_2}=24/4=6$.
\[
  \operatorname{RSS}_{R_1}=(2-3)^2+(4-3)^2=2,\quad
  \operatorname{RSS}_{R_2}=(3-6)^2+(3-6)^2+(8-6)^2+(10-6)^2=9+9+4+16=38,
\]
so $\operatorname{RSS}(2.5)=2+38=\mathbf{40}$.

\emph{Cutpoint $s=4.5$:} $R_1=\{1,2,3,4\}$ with $y$-values $\{2,4,3,3\}$, mean
$\hat{y}_{R_1}=12/4=3$; $R_2=\{5,6\}$ with $y$-values $\{8,10\}$, mean $\hat{y}_{R_2}=9$.
\[
  \operatorname{RSS}_{R_1}=1+1+0+0=2,\qquad
  \operatorname{RSS}_{R_2}=(8-9)^2+(10-9)^2=2,
\]
so $\operatorname{RSS}(4.5)=2+2=\mathbf{4}$.

Since $4<40$, recursive binary splitting is \emph{greedy} and picks the cutpoint with the
smallest RSS: $\mathbf{s=4.5}$. (It separates the four low-sales products, mean $3$, from the
two high-sales products, mean $9$.)

\medskip
\textbf{(b) Prediction for $x=5$.}
$x=5\ge 4.5$, so the new product falls into $R_2$ and the tree predicts
$\hat{y}=\hat{y}_{R_2}=\mathbf{9}$ (i.e.\ $9{,}000$ units per week).

\medskip
\textbf{(c) Node impurity.}
$\hat{p}_{\text{Yes}}=6/8=0.75$, $\hat{p}_{\text{No}}=2/8=0.25$.
\[
  G=0.75\cdot 0.25+0.25\cdot 0.75=0.1875+0.1875=\mathbf{0.375}.
\]
\[
  D=-\bigl(0.75\ln 0.75+0.25\ln 0.25\bigr)
   =-\bigl(0.75\cdot(-0.2877)+0.25\cdot(-1.3863)\bigr)
   =0.2158+0.3466=\mathbf{0.562}.
\]
The node predicts the majority class, \textbf{Yes} (with estimated probability $0.75$).
Both $G$ and $D$ are small when the node is pure; here the node is moderately impure.

\medskip
\textbf{(d) Role of $\alpha$.}
$\alpha$ is the price per terminal node: it trades off the fit of the subtree (first term)
against its complexity $|T|$.
For $\alpha=0$ the criterion equals the training RSS and the full tree $T_0$ is optimal.
As $\alpha$ increases, additional leaves must ``pay for themselves''; branches whose RSS
reduction is smaller than $\alpha$ are pruned, so the optimal subtree becomes smaller ---
one obtains a \emph{nested sequence} of subtrees as $\alpha$ grows (weakest-link pruning).
In practice $\alpha$ is selected by \emph{cross-validation} (choose the $\alpha$ with the
smallest CV error, then refit on all data), which balances bias against variance instead of
optimising training fit.
\end{solutionbox}
\fi

% ============================================================
\problem{4: Ensembles --- bagging, random forests, boosting}{16}

\begin{enumerate}[label=(\alph*),leftmargin=7mm,itemsep=2.5mm]
  \item \pp{6} Let $\hat{f}_1(x),\dots,\hat{f}_B(x)$ be $B$ identically distributed (but not
    independent) tree predictions at a fixed point $x$, each with variance $\sigma^2$ and with
    pairwise correlation $\rho$. The variance of the bagged average is
    \[
      \operatorname{Var}\Bigl(\tfrac{1}{B}\sum_{b=1}^{B}\hat{f}_b(x)\Bigr)
      =\rho\,\sigma^{2}+\frac{1-\rho}{B}\,\sigma^{2}.
    \]
    (i) Explain briefly why this formula shows that bagging reduces variance, and which term
    limits the reduction. (ii) Evaluate the variance for $B=100$, $\rho=0.6$ and $\sigma^2=4$,
    and compare it with the variance of a single tree. (iii) What value does the variance
    approach as $B\to\infty$?
  \item \pp{4} How do \emph{random forests} modify bagging at each split, and what is the
    recommended subset size $m$ for classification as a function of $p$ (evaluate for $p=25$)?
    Explain, with reference to the formula in (a), why this modification improves on bagging.
  \item \pp{4} Name the three tuning parameters of \emph{boosting} for regression trees and
    state in one sentence each what they control. What is the ``slow learning'' idea behind
    boosting?
  \item \pp{2} What does the \emph{out-of-bag} (OOB) error of a bagged model estimate, and why
    is it available essentially for free?
\end{enumerate}

\ifdefined\withsolutions
\begin{solutionbox}
\textbf{(a) Variance of the bagged average.}
(i) The second term $\frac{1-\rho}{B}\sigma^2$ is driven to $0$ by averaging over many trees
($B$ large), so averaging removes the ``independent part'' of the trees' variability ---
this is the variance reduction of bagging. The first term $\rho\sigma^2$ does \emph{not}
depend on $B$: because the bootstrap trees are grown on overlapping samples of the same data
(and share the same strong predictors), they are positively correlated, and this correlated
part of the variance cannot be averaged away. It limits the reduction.
(ii) For $B=100$, $\rho=0.6$, $\sigma^2=4$:
\[
  \operatorname{Var}=0.6\cdot 4+\frac{1-0.6}{100}\cdot 4
  =2.4+\frac{0.4\cdot 4}{100}=2.4+0.016=\mathbf{2.416},
\]
compared with $\sigma^2=4$ for a single tree --- a reduction of about $40\%$.
(iii) As $B\to\infty$ the variance approaches the floor $\rho\sigma^2=0.6\cdot 4=\mathbf{2.4}$.
(Increasing $B$ further does not overfit; it only stabilises the average.)

\medskip
\textbf{(b) Random forests.}
At each split, only a \emph{random subset} of $m$ of the $p$ predictors is considered as
splitting candidates; a fresh subset is drawn at every split. Recommended for classification:
$m\approx\sqrt{p}$, so for $p=25$: $m=\sqrt{25}=\mathbf{5}$.
Why it helps: with bagging, a dominant predictor is used in the top split of almost every tree,
so the trees look alike and $\rho$ is high. Excluding it from $\approx (p-m)/p$ of the split
searches \emph{decorrelates} the trees, i.e.\ lowers $\rho$ --- and by (a) this lowers the
irreducible floor $\rho\sigma^2$ of the ensemble variance, which mere averaging cannot touch.

\medskip
\textbf{(c) Boosting.}
\begin{itemize}[nosep,leftmargin=5mm]
  \item \emph{Number of trees $B$:} how many small trees are added sequentially; unlike
    bagging, boosting \emph{can} overfit if $B$ is too large, so $B$ is chosen by
    cross-validation.
  \item \emph{Shrinkage $\lambda$} (learning rate, e.g.\ $0.01$ or $0.001$): scales the
    contribution of each new tree and thus the speed of learning; smaller $\lambda$ usually
    needs larger $B$.
  \item \emph{Number of splits $d$ per tree} (interaction depth): the complexity of each tree
    and the interaction order of the model; often even $d=1$ (stumps, additive model) works
    well.
\end{itemize}
\emph{Slow learning:} each tree is fitted to the current \emph{residuals}, and only a small
fraction $\lambda$ of it is added to the model; the model thus improves slowly and exactly in
the regions where it still performs poorly, instead of fitting the data hard in one step ---
which tends to generalise better.

\medskip
\textbf{(d) OOB error.}
Each bootstrap sample leaves out on average about $1/3$ of the observations
(``out-of-bag''). Predicting each observation only with the trees that did \emph{not} use it,
and aggregating these predictions, yields the OOB error --- a valid estimate of the
\emph{test error} of the bagged model. It is free because it reuses the bootstrap side
product; no separate validation set or cross-validation loop is required.
\end{solutionbox}
\fi

% ============================================================
\problem{5: Neural networks by hand}{20}

Consider a tiny feed-forward network for a quantitative response with input
$x=(x_1,x_2)=(1,2)$, one hidden layer with two ReLU units
($g(z)=\max\{0,z\}$) and a linear output unit:
\[
  A_1=g\bigl(-1+1\,x_1+2\,x_2\bigr),\qquad
  A_2=g\bigl(-3-2\,x_1+1\,x_2\bigr),\qquad
  f(x)=1+2\,A_1+4\,A_2 .
\]

\begin{enumerate}[label=(\alph*),leftmargin=7mm,itemsep=2.5mm]
  \item \pp{8} Perform the forward pass step by step: compute both pre-activations, both
    activations $A_1,A_2$ and the network output $f(x)$. Comment briefly on what the ReLU did
    to the second unit and why such non-linear activations are essential.
  \item \pp{6} A classification network with $K=3$ classes produces the logits (output-layer
    pre-activations) $z=(z_1,z_2,z_3)=(1.0,\ 0.0,\ -1.0)$ for one observation whose true class
    is class~1. Compute all three softmax probabilities and the cross-entropy loss
    $-\ln \hat{p}_{\text{true}}$.
    \emph{(Use $e^{1}=2.718$, $e^{0}=1.000$, $e^{-1}=0.368$ and $\ln 4.086=1.408$.)}
  \item \pp{4} Count the parameters of a fully connected network with $p=4$ inputs, one hidden
    layer with $k=3$ units and an output layer with $K=2$ units, where every hidden and output
    unit has a bias. Show the count separately for the two weight layers.
  \item \pp{2} A convolutional layer receives a $32\times 32$ image and uses filters of size
    $F=5$ with padding $P=2$ and stride $S=1$. Compute the spatial size of the output feature
    map from $\bigl(W-F+2P\bigr)/S+1$.
\end{enumerate}

\ifdefined\withsolutions
\begin{solutionbox}
\textbf{(a) Forward pass.}
Pre-activations:
\[
  z_1=-1+1\cdot 1+2\cdot 2=-1+1+4=4,\qquad
  z_2=-3-2\cdot 1+1\cdot 2=-3-2+2=-3 .
\]
ReLU activations:
\[
  A_1=g(4)=\max\{0,4\}=4,\qquad A_2=g(-3)=\max\{0,-3\}=0 .
\]
Output:
\[
  f(x)=1+2\cdot 4+4\cdot 0=1+8+0=\mathbf{9}.
\]
The ReLU \emph{switched the second unit off}: its negative pre-activation is clipped to $0$,
so unit 2 contributes nothing for this input (but would for other inputs). This thresholding
is the source of non-linearity: without a non-linear $g$, the network would collapse into a
single linear function of $x$, no matter how many layers it has.

\medskip
\textbf{(b) Softmax and cross-entropy.}
Exponentials: $e^{1.0}=2.718$, $e^{0.0}=1.000$, $e^{-1.0}=0.368$; their sum is
$2.718+1.000+0.368=4.086$. Softmax probabilities:
\[
  \hat{p}_1=\frac{2.718}{4.086}=0.665,\qquad
  \hat{p}_2=\frac{1.000}{4.086}=0.245,\qquad
  \hat{p}_3=\frac{0.368}{4.086}=0.090,
\]
which sum to $1.000$. Cross-entropy loss for true class 1:
\[
  -\ln\hat{p}_1=-\ln\frac{e^{1}}{4.086}=-\bigl(1-\ln 4.086\bigr)=1.408-1=\mathbf{0.408}.
\]
(Equivalently $-\ln 0.665\approx 0.408$; a perfect prediction $\hat{p}_1=1$ would give loss 0.)

\medskip
\textbf{(c) Parameter count for $4\to 3\to 2$.}
Input $\to$ hidden: $4\cdot 3=12$ weights $+\,3$ biases $=15$.
Hidden $\to$ output: $3\cdot 2=6$ weights $+\,2$ biases $=8$.
Total: $15+8=\mathbf{23}$ parameters.

\medskip
\textbf{(d) Convolution output size.}
\[
  \frac{W-F+2P}{S}+1=\frac{32-5+2\cdot 2}{1}+1=\frac{31}{1}+1=32,
\]
so the output feature map is $\mathbf{32\times 32}$ --- padding $P=2$ with $F=5$, $S=1$ is
exactly the ``same'' convolution that preserves the spatial size.
\end{solutionbox}
\fi

% ============================================================
\problem{6: Multiple testing}{20}

\begin{enumerate}[label=(\alph*),leftmargin=7mm,itemsep=2.5mm]
  \item \pp{4} A lab performs $m=20$ \emph{independent} hypothesis tests, each at level
    $\alpha=0.05$, and suppose all 20 null hypotheses are true. Define the family-wise error
    rate (FWER) and compute it for this situation
    \emph{(use $(0.95)^{20}=0.3585$)}. Interpret the result in one sentence.
  \item \pp{7} A research group screens $m=6$ candidate biomarkers; the ordered $p$-values are
    \begin{center}
    \begin{tabular}{lcccccc}
    \toprule
    $j$ & 1 & 2 & 3 & 4 & 5 & 6\\
    \midrule
    $p_{(j)}$ & 0.001 & 0.004 & 0.011 & 0.030 & 0.045 & 0.230\\
    \bottomrule
    \end{tabular}
    \end{center}
    Control the FWER at $\alpha=0.05$ (i) with \emph{Bonferroni} and (ii) with the
    \emph{Holm step-down} procedure. Give the relevant thresholds, state for each method
    exactly which hypotheses are rejected, and explain why Holm can never reject fewer
    hypotheses than Bonferroni.
  \item \pp{6} Apply the \emph{Benjamini--Hochberg} (BH) procedure with FDR level $q=0.05$ to
    the same six $p$-values: give the thresholds $jq/m$, determine which hypotheses are
    rejected, and state precisely what ``FDR control at $q=0.05$'' guarantees for the
    resulting set of rejections.
  \item \pp{3} A genomics group screens $20{,}000$ genes in an \emph{exploratory} study; the
    most promising hits will be re-tested in an independent follow-up experiment. Should they
    control the FWER or the FDR? Justify your recommendation.
\end{enumerate}

\ifdefined\withsolutions
\begin{solutionbox}
\textbf{(a) FWER for $m=20$ independent tests.}
$\text{FWER}=\Pr(\text{at least one Type I error})=\Pr(V\ge 1)$, where $V$ is the number of
true nulls that are (falsely) rejected. With independence and all nulls true:
\[
  \text{FWER}=1-\Pr(\text{no false rejection})=1-(1-0.05)^{20}=1-(0.95)^{20}
  =1-0.3585=\mathbf{0.6415}.
\]
Interpretation: even though each single test is performed at the $5\%$ level, there is a
$64\%$ chance of at least one false discovery among the 20 tests --- testing many hypotheses
without adjustment almost guarantees spurious findings.

\medskip
\textbf{(b) Bonferroni and Holm at $\alpha=0.05$.}
\emph{Bonferroni:} reject $H_{(j)}$ iff $p_{(j)}\le\alpha/m=0.05/6=0.00833$.
$p_{(1)}=0.001\le 0.00833$ \checkmark, $p_{(2)}=0.004\le 0.00833$ \checkmark,
$p_{(3)}=0.011>0.00833$ $\times$.
$\Rightarrow$ Bonferroni rejects $H_{(1)}$ and $H_{(2)}$: \textbf{2 rejections}.

\emph{Holm:} compare $p_{(j)}$ with $\alpha/(m-j+1)$ step by step and stop at the first failure:
\begin{center}
\begin{tabular}{lcccccc}
\toprule
$j$ & 1 & 2 & 3 & 4 & 5 & 6\\
\midrule
$p_{(j)}$ & 0.001 & 0.004 & 0.011 & 0.030 & 0.045 & 0.230\\
$\alpha/(m-j+1)$ & 0.00833 & 0.0100 & 0.0125 & 0.0167 & 0.0250 & 0.0500\\
reject? & \checkmark & \checkmark & \checkmark & stop & --- & ---\\
\bottomrule
\end{tabular}
\end{center}
$0.001\le 0.00833$, $0.004\le 0.0100$, $0.011\le 0.0125$, but $0.030>0.0167$: stop.
$\Rightarrow$ Holm rejects $H_{(1)},H_{(2)},H_{(3)}$: \textbf{3 rejections}.

\emph{Why Holm $\supseteq$ Bonferroni:} Holm's thresholds $\alpha/(m-j+1)$ are $\ge\alpha/m$
for every $j$ (with equality only at $j=1$), so every $p$-value that passes Bonferroni also
passes Holm's less stringent later hurdles --- Holm is uniformly more powerful, yet still
controls the FWER at $\alpha$ without any independence assumption.

\medskip
\textbf{(c) Benjamini--Hochberg at $q=0.05$.}
Thresholds $jq/m=j\cdot 0.05/6$:
\begin{center}
\begin{tabular}{lcccccc}
\toprule
$j$ & 1 & 2 & 3 & 4 & 5 & 6\\
\midrule
$p_{(j)}$ & 0.001 & 0.004 & 0.011 & 0.030 & 0.045 & 0.230\\
$jq/m$ & 0.00833 & 0.0167 & 0.0250 & 0.0333 & 0.0417 & 0.0500\\
$p_{(j)}\le jq/m$? & yes & yes & yes & \textbf{yes} & no & no\\
\bottomrule
\end{tabular}
\end{center}
$L=$ largest $j$ with $p_{(j)}\le jq/m$: $j=4$ (since $0.030\le 0.0333$, while
$0.045>0.0417$ and $0.230>0.0500$). BH rejects \emph{all} $H_{(j)}$ with $j\le L$:
$H_{(1)},\dots,H_{(4)}$: \textbf{4 rejections}.
Guarantee: $\text{FDR}=E\bigl[V/\max(R,1)\bigr]\le q=0.05$ --- \emph{on average} at most
$5\%$ of the rejected hypotheses are false discoveries. It does \emph{not} guarantee that at
most $5\%$ of these particular 4 rejections are false, nor does it bound the probability of
making any false rejection (that would be FWER control).

\medskip
\textbf{(d) Exploratory genomics screening.}
Control the \textbf{FDR}. With $m=20{,}000$ tests, FWER control (e.g.\ Bonferroni at
$0.05/20{,}000=2.5\cdot 10^{-6}$) guards against even a single false positive and therefore
has extremely low power --- almost no true signals would survive. In an exploratory screen
the cost of a few false leads is low, because every hit is validated in an independent
follow-up experiment; what matters is that the \emph{list} of discoveries is mostly reliable.
FDR control at, say, $q=0.05$--$0.2$ delivers exactly that: a rich list of candidates of
which on average at most a known fraction is false.
\end{solutionbox}
\fi

% ============================================================
\problem{7: Cumulative essentials (Chapters 2--6)}{20}

\begin{enumerate}[label=(\alph*),leftmargin=7mm,itemsep=2.5mm]
  \item \pp{4} True or false --- justify briefly: ``A more flexible statistical learning method
    always yields a lower \emph{test} error than a less flexible one.''
  \item \pp{4} A logistic regression for loan default,
    $\log\bigl(\frac{p(x)}{1-p(x)}\bigr)=\beta_0+\beta_1 x$ with $x=$ number of missed
    payments, yields $\hat{\beta}_1=0.7$. Interpret the effect of \emph{one additional missed
    payment} on the \emph{odds} of default \emph{(use $e^{0.7}=2.014$)}. Why can the effect on
    the \emph{probability} of default not be stated as a single number?
  \item \pp{4} Give one reason to prefer $k$-fold cross-validation with $k=5$ or $10$ over
    \emph{LOOCV}, and one reason to prefer it over a \emph{single validation split}.
  \item \pp{2} You have $p=2{,}000$ predictors but expect only a handful to be truly related to
    the response. Ridge or lasso --- which penalty is more appropriate, and why?
  \item \pp{6} Match each scenario to exactly one method from
    \{\emph{linear regression}, \emph{lasso}, \emph{logistic regression},
    \emph{random forest}\} (each method used exactly once) and justify each choice in one
    sentence:
    \begin{enumerate}[label=(\roman*),nosep,topsep=1mm]
      \item Estimate how much an additional square metre adds to the price of an apartment,
        with a confidence interval, from $n=2{,}000$ sales.
      \item Predict a continuous drug response from $p=5{,}000$ gene-expression measurements
        for $n=100$ patients; clinicians want a short list of relevant genes.
      \item Predict whether a loan applicant will default (yes/no) and report to the regulator
        how each applicant characteristic changes the odds of default.
      \item Predict machine failures from hundreds of sensor variables with strong non-linear
        effects and interactions; only predictive accuracy matters.
    \end{enumerate}
\end{enumerate}

\ifdefined\withsolutions
\begin{solutionbox}
\textbf{(a) False.}
Expected test error $=$ variance $+$ bias$^2$ $+$ irreducible error. More flexibility reduces
bias but \emph{increases} variance; once the variance increase outweighs the bias reduction,
the test error rises again --- the characteristic U-shape. So a more flexible method wins only
up to the sweet spot; e.g.\ when the true $f$ is nearly linear or $n$ is small, a flexible
method overfits and a simple linear fit has \emph{lower} test error. (Only the \emph{training}
error always decreases with flexibility.)

\medskip
\textbf{(b) Odds ratio.}
One additional missed payment increases the log-odds by $\hat{\beta}_1=0.7$, i.e.\ it
\emph{multiplies the odds} of default by $e^{0.7}=2.014$ --- the odds roughly double,
regardless of the starting value of $x$.
The effect on the \emph{probability} depends on where the observation sits on the S-shaped
logistic curve: the same change in log-odds shifts $p(x)$ a lot near $p\approx 0.5$ but very
little when $p$ is close to $0$ or $1$ --- hence no single number describes it.

\medskip
\textbf{(c) $k$-fold CV with $k=5$ or $10$.}
\emph{vs.\ LOOCV:} the $n$ LOOCV fits are trained on almost identical data sets, so their
errors are highly positively correlated; averaging strongly correlated quantities reduces
variance little, so the LOOCV estimate has \emph{higher variance} than 5- or 10-fold CV
(it is also computationally expensive: $n$ fits instead of $k$, absent shortcut formulas).
\emph{vs.\ a single validation split:} the validation estimate depends strongly on the random
split and overestimates the test error because only part of the data is used for training
(\emph{bias}); $k$-fold CV trains on $(k-1)/k$ of the data and averages over $k$ folds, giving
a less biased and more stable estimate.

\medskip
\textbf{(d) Lasso.}
The $\ell_1$ penalty sets many coefficients \emph{exactly to zero} and thus performs variable
selection --- ideal when the truth is sparse (a handful of signals among 2{,}000 predictors),
and it yields an interpretable small model. Ridge shrinks coefficients but keeps all
$p=2{,}000$ predictors in the model and tends to lose when most true coefficients are zero.

\medskip
\textbf{(e) Matching.}
\begin{itemize}[nosep,leftmargin=5mm]
  \item (i) \textbf{Linear regression}: quantitative response, goal is \emph{inference}
    (effect size with confidence interval) --- exactly what OLS with standard errors provides.
  \item (ii) \textbf{Lasso}: $p\gg n$ makes OLS infeasible; the $\ell_1$ penalty regularises
    and selects a short list of genes.
  \item (iii) \textbf{Logistic regression}: binary response with interpretable coefficients
    ($e^{\hat\beta_j}$ = odds ratios) --- suitable for regulatory reporting.
  \item (iv) \textbf{Random forest}: captures non-linearities and interactions automatically,
    is accurate and robust with many mixed predictors; interpretability is not required.
\end{itemize}
\end{solutionbox}
\fi

\vspace{6mm}
\begin{center}
\rule{0.35\textwidth}{0.4pt}\\[1mm]
\textit{End of the examination --- good luck!}
\end{center}

\end{document}
"
% ============================================================
\documentclass[11pt,a4paper]{article}
\usepackage[utf8]{inputenc}
\usepackage[T1]{fontenc}
\usepackage[english]{babel}
\usepackage[margin=2.2cm]{geometry}
\usepackage{amsmath,amssymb}
\usepackage{booktabs,array}
\usepackage{enumitem}
\usepackage{xcolor}
\usepackage{tcolorbox}
\tcbuselibrary{skins,breakable}

\setlength{\parindent}{0pt}

% Teal solution box (only shown when \withsolutions is defined)
\newtcolorbox{solutionbox}{enhanced,breakable,colback=teal!5,colframe=teal!55!black,
  boxrule=0pt,leftrule=2.5pt,arc=2pt,fonttitle=\bfseries,title=Solution,
  left=2mm,right=2mm,top=1mm,bottom=1.5mm}

\newcommand{\problem}[2]{\par\vspace{7mm}\noindent{\large\textbf{Problem #1}}%
  \hfill{\large\textbf{(#2 points)}}\par\nopagebreak\vspace{0.6mm}%
  \noindent\rule{\textwidth}{0.6pt}\par\nopagebreak\vspace{2.5mm}}
\newcommand{\pp}[1]{\textbf{[#1~P]}\enspace}

\begin{document}

% ------------------------------------------------------------
% Header block
% ------------------------------------------------------------
\begin{center}
  {\Large\textbf{HSBI --- Bielefeld University of Applied Sciences and Arts}}\\[1.2mm]
  {\large Quantitative Research Methods}\\[1mm]
  {\normalsize Summer Semester 2026 \quad$\cdot$\quad Examiner: Prof.\ Dr.\ Christoph Weisser}\\[4.5mm]
  {\LARGE\textbf{Final Mock Examination (Lectures 1--12)}}\\[1.2mm]
  {\normalsize Comprehensive coverage of ISLP Chapters 2--8, 10 and 13,\\
   with emphasis on Chapters 7 (splines/GAMs), 8 (trees), 10 (deep learning) and 13 (multiple testing)}%
  \ifdefined\withsolutions\\[2.2mm]{\large\textcolor{teal!55!black}{\textbf{--- Version with detailed solutions ---}}}\fi
\end{center}

\vspace{2.5mm}
\begin{center}
\begin{tabular}{@{}ll@{}}
  \textbf{Duration:} & \textbf{120 minutes}\\
  \textbf{Total credit:} & \textbf{120 points}\\
  \textbf{Allowed aids:} & non-programmable calculator; one handwritten A4 sheet of notes (both sides)\\
\end{tabular}
\end{center}

\vspace{3mm}
Name: \rule[-1mm]{0.38\textwidth}{0.4pt}\hfill Matriculation no.: \rule[-1mm]{0.24\textwidth}{0.4pt}

\vspace{4.5mm}
\begin{center}
\renewcommand{\arraystretch}{1.35}
\begin{tabular}{|l|c|c|c|c|c|c|c|c|}
\hline
\textbf{Problem} & \textbf{P1} & \textbf{P2} & \textbf{P3} & \textbf{P4} & \textbf{P5} & \textbf{P6} & \textbf{P7} & \textbf{Total}\\\hline
Maximum points & 16 & 8 & 20 & 16 & 20 & 20 & 20 & \textbf{120}\\\hline
Achieved points & \hspace{8mm} & \hspace{8mm} & \hspace{8mm} & \hspace{8mm} & \hspace{8mm} & \hspace{8mm} & \hspace{8mm} & \hspace{10mm}\\\hline
\end{tabular}
\end{center}

\vspace{2.5mm}
\textbf{Instructions}
\begin{itemize}[nosep,leftmargin=5.5mm]
  \item Justify all answers. Numerical results without a derivation or a brief reasoning receive no credit.
  \item Round final numerical results to 3 decimal places unless stated otherwise. Small deviations caused by carrying rounded intermediate results are accepted.
  \item All numerical constants you need (powers of $e$, logarithms, $(0.95)^{20}$, $t$-quantiles) are provided in the problems.
  \item The number of points per problem roughly equals the recommended working time in minutes.
\end{itemize}
\vspace{1mm}\noindent\rule{\textwidth}{0.6pt}

% ============================================================
\problem{1: Moving beyond linearity --- splines}{16}

\begin{enumerate}[label=(\alph*),leftmargin=7mm,itemsep=2.5mm]
  \item \pp{4} A \emph{cubic spline} is a piecewise cubic polynomial that is continuous and has
    continuous first and second derivatives at each knot. How many degrees of freedom (free
    parameters, counting the intercept) does a cubic spline with $K$ interior knots use?
    Justify your answer by counting the parameters of the polynomial pieces and the constraints
    at the knots, and evaluate the result for $K=4$.
  \item \pp{4} How many degrees of freedom does a \emph{natural} cubic spline with the same $K$
    knots use (again evaluate for $K=4$)? Which additional constraint does a natural spline
    impose, in which region of the predictor space does it act, and why is this constraint
    useful in practice?
  \item \pp{4} A \emph{smoothing spline} is the function $g$ that minimises
    \[
      \sum_{i=1}^{n}\bigl(y_i-g(x_i)\bigr)^2+\lambda\int g''(t)^2\,dt .
    \]
    Explain the role of the two terms, and describe the behaviour of the solution $\hat{g}$ in
    the two limiting cases $\lambda\to 0$ and $\lambda\to\infty$.
  \item \pp{4} Write down the \emph{truncated-power basis} representation of a cubic spline with
    a single knot at $\xi$, define the truncated-power basis function explicitly, and explain in
    one sentence why adding this basis function does not destroy the smoothness of the fit at
    $\xi$.
\end{enumerate}

\ifdefined\withsolutions
\begin{solutionbox}
\textbf{(a) Degrees of freedom of a cubic spline.}
$K$ interior knots split the range of $X$ into $K+1$ intervals, each carrying a cubic
polynomial with $4$ coefficients: $4(K+1)=4K+4$ parameters. At each knot, continuity of
$g$, $g'$ and $g''$ imposes $3$ constraints, i.e.\ $3K$ constraints in total. Hence
\[
  \mathrm{df}=4(K+1)-3K=K+4 .
\]
For $K=4$: $\mathrm{df}=4+4=\mathbf{8}$ free parameters (equivalently: intercept, $x$, $x^2$,
$x^3$ plus one truncated-power basis function per knot $=4+K$).

\medskip
\textbf{(b) Natural cubic spline.}
A natural spline additionally requires the function to be \emph{linear beyond the boundary
knots}, i.e.\ in the two outer regions $X<\xi_1$ and $X>\xi_K$. Forcing the quadratic and cubic
terms to vanish there removes $2$ parameters per boundary region, i.e.\ $4$ in total:
\[
  \mathrm{df}=(K+4)-4=K,\qquad K=4:\ \mathrm{df}=\mathbf{4}.
\]
Usefulness: polynomials and unconstrained splines are notoriously erratic near the boundaries,
where data are sparse; the linearity constraint stabilises the fit there
(\emph{lower variance} of predictions at and beyond the boundary), at the price of a small
amount of additional bias.

\medskip
\textbf{(c) Smoothing spline.}
The first term (RSS) rewards fidelity to the data; the second term penalises roughness, since
$g''$ measures the curvature (wiggliness) of $g$, and $\lambda\ge 0$ is the tuning parameter
that trades the two off.
\emph{$\lambda\to 0$:} the penalty disappears; $\hat{g}$ can be made arbitrarily wiggly and
\emph{interpolates} the training observations exactly (zero training RSS, severe overfitting;
effective df $\to n$).
\emph{$\lambda\to\infty$:} any curvature is infinitely expensive, so $g''\equiv 0$ is enforced;
$\hat{g}$ becomes a straight line and equals the \emph{least-squares regression line}
(effective df $\to 2$). Between the extremes, $\lambda$ (chosen e.g.\ by LOOCV) controls the
effective degrees of freedom, which decrease smoothly from $n$ to $2$.

\medskip
\textbf{(d) Truncated-power basis with one knot $\xi$.}
\[
  f(x)=\beta_0+\beta_1x+\beta_2x^2+\beta_3x^3+\beta_4\,(x-\xi)^3_+,
  \qquad
  (x-\xi)^3_+=\begin{cases}(x-\xi)^3, & x>\xi\\[0.5mm] 0, & x\le\xi\end{cases}
\]
The function $(x-\xi)^3_+$ and its first and second derivatives are all $0$ at $x=\xi$, so
adding it changes only the third derivative at the knot; the fit remains continuous with
continuous first and second derivatives --- exactly the defining smoothness of a cubic spline.
\end{solutionbox}
\fi

% ============================================================
\problem{2: Generalised additive models}{8}

\begin{enumerate}[label=(\alph*),leftmargin=7mm,itemsep=2.5mm]
  \item \pp{3} Write down the general form of a \emph{generalised additive model} (GAM) for a
    quantitative response $Y$ and predictors $X_1,\dots,X_p$, and explain in one sentence how it
    extends multiple linear regression.
  \item \pp{3} State one important \emph{advantage} of GAMs that concerns the interpretation of
    the fitted functions $f_j$, and the main structural \emph{limitation} of GAMs --- including
    how this limitation can be repaired.
  \item \pp{2} Describe in one or two sentences how the \emph{backfitting} algorithm fits a GAM.
\end{enumerate}

\ifdefined\withsolutions
\begin{solutionbox}
\textbf{(a) Model form.}
\[
  Y=\beta_0+f_1(X_1)+f_2(X_2)+\dots+f_p(X_p)+\varepsilon .
\]
Multiple linear regression is the special case $f_j(X_j)=\beta_jX_j$; a GAM replaces each linear
component by a smooth non-linear function $f_j$ (e.g.\ a natural or smoothing spline, or local
regression), while keeping the \emph{additive} structure.

\medskip
\textbf{(b) Advantage and limitation.}
\emph{Advantage (interpretability):} because the model is additive, the contribution of each
$X_j$ to $Y$ is captured by its own function $f_j$, which can be examined and plotted
individually while the other variables are held fixed --- non-linear effects remain as easy to
read as in a coefficient table, one panel per predictor.
\emph{Limitation:} additivity excludes \emph{interactions}: the effect of $X_j$ cannot depend
on the value of $X_k$. Repair: add interaction terms manually, e.g.\ a product term
$X_j\times X_k$ or a low-dimensional smooth surface $f_{jk}(X_j,X_k)$ as an extra component.

\medskip
\textbf{(c) Backfitting.}
Cycle through the predictors and repeatedly re-fit each $f_j$ (by a smoother, e.g.\ a spline)
to the \emph{partial residuals} $r_i=y_i-\beta_0-\sum_{k\ne j}f_k(x_{ik})$, holding all other
fitted functions fixed; iterate until the $f_j$ stop changing.
\end{solutionbox}
\fi

% ============================================================
\problem{3: Regression and classification trees by hand}{20}

An online platform records for $n=6$ products the advertising budget $x$ (in 1{,}000 EUR) and
the weekly sales $y$ (in 1{,}000 units):

\begin{center}
\begin{tabular}{lcccccc}
\toprule
Product $i$ & 1 & 2 & 3 & 4 & 5 & 6\\
\midrule
Budget $x_i$ & 1 & 2 & 3 & 4 & 5 & 6\\
Sales $y_i$ & 2 & 4 & 3 & 3 & 8 & 10\\
\bottomrule
\end{tabular}
\end{center}

A regression tree with a \emph{single} split at cutpoint $s$ partitions the data into
$R_1=\{i: x_i<s\}$ and $R_2=\{i: x_i\ge s\}$ and predicts the region mean
$\hat{y}_{R_m}$ in each region.

\begin{enumerate}[label=(\alph*),leftmargin=7mm,itemsep=2.5mm]
  \item \pp{8} For each of the two candidate cutpoints $s=2.5$ and $s=4.5$, determine the two
    regions, the region means and the resulting total residual sum of squares
    $\operatorname{RSS}(s)=\sum_{m=1}^{2}\sum_{i\in R_m}(y_i-\hat{y}_{R_m})^2$.
    Which cutpoint does recursive binary splitting choose, and why?
  \item \pp{2} Using the chosen split, predict the weekly sales for a new product with
    advertising budget $x=5$.
  \item \pp{6} A \emph{classification} tree for a different task has a terminal node containing
    8 training observations: 6 of class ``Yes'' and 2 of class ``No''. Compute the class
    proportions, the Gini index $G=\sum_{k}\hat{p}_{k}(1-\hat{p}_{k})$ and the entropy
    $D=-\sum_{k}\hat{p}_{k}\ln\hat{p}_{k}$ of this node
    \emph{(use $\ln 0.75=-0.2877$ and $\ln 0.25=-1.3863$)}.
    Which class does the node predict?
  \item \pp{4} In \emph{cost-complexity pruning} one minimises, for a subtree
    $T\subseteq T_0$ of the large tree $T_0$,
    \[
      \sum_{m=1}^{|T|}\ \sum_{i:\,x_i\in R_m}\bigl(y_i-\hat{y}_{R_m}\bigr)^2+\alpha\,|T| .
    \]
    Explain the role of the tuning parameter $\alpha$: what happens for $\alpha=0$, what happens
    as $\alpha$ increases, and how is $\alpha$ chosen in practice?
\end{enumerate}

\ifdefined\withsolutions
\begin{solutionbox}
\textbf{(a) Evaluating the two cutpoints.}
Root node for reference: $\bar{y}=30/6=5$.

\emph{Cutpoint $s=2.5$:} $R_1=\{1,2\}$ with $y$-values $\{2,4\}$, mean $\hat{y}_{R_1}=3$;
$R_2=\{3,4,5,6\}$ with $y$-values $\{3,3,8,10\}$, mean $\hat{y}_{R_2}=24/4=6$.
\[
  \operatorname{RSS}_{R_1}=(2-3)^2+(4-3)^2=2,\quad
  \operatorname{RSS}_{R_2}=(3-6)^2+(3-6)^2+(8-6)^2+(10-6)^2=9+9+4+16=38,
\]
so $\operatorname{RSS}(2.5)=2+38=\mathbf{40}$.

\emph{Cutpoint $s=4.5$:} $R_1=\{1,2,3,4\}$ with $y$-values $\{2,4,3,3\}$, mean
$\hat{y}_{R_1}=12/4=3$; $R_2=\{5,6\}$ with $y$-values $\{8,10\}$, mean $\hat{y}_{R_2}=9$.
\[
  \operatorname{RSS}_{R_1}=1+1+0+0=2,\qquad
  \operatorname{RSS}_{R_2}=(8-9)^2+(10-9)^2=2,
\]
so $\operatorname{RSS}(4.5)=2+2=\mathbf{4}$.

Since $4<40$, recursive binary splitting is \emph{greedy} and picks the cutpoint with the
smallest RSS: $\mathbf{s=4.5}$. (It separates the four low-sales products, mean $3$, from the
two high-sales products, mean $9$.)

\medskip
\textbf{(b) Prediction for $x=5$.}
$x=5\ge 4.5$, so the new product falls into $R_2$ and the tree predicts
$\hat{y}=\hat{y}_{R_2}=\mathbf{9}$ (i.e.\ $9{,}000$ units per week).

\medskip
\textbf{(c) Node impurity.}
$\hat{p}_{\text{Yes}}=6/8=0.75$, $\hat{p}_{\text{No}}=2/8=0.25$.
\[
  G=0.75\cdot 0.25+0.25\cdot 0.75=0.1875+0.1875=\mathbf{0.375}.
\]
\[
  D=-\bigl(0.75\ln 0.75+0.25\ln 0.25\bigr)
   =-\bigl(0.75\cdot(-0.2877)+0.25\cdot(-1.3863)\bigr)
   =0.2158+0.3466=\mathbf{0.562}.
\]
The node predicts the majority class, \textbf{Yes} (with estimated probability $0.75$).
Both $G$ and $D$ are small when the node is pure; here the node is moderately impure.

\medskip
\textbf{(d) Role of $\alpha$.}
$\alpha$ is the price per terminal node: it trades off the fit of the subtree (first term)
against its complexity $|T|$.
For $\alpha=0$ the criterion equals the training RSS and the full tree $T_0$ is optimal.
As $\alpha$ increases, additional leaves must ``pay for themselves''; branches whose RSS
reduction is smaller than $\alpha$ are pruned, so the optimal subtree becomes smaller ---
one obtains a \emph{nested sequence} of subtrees as $\alpha$ grows (weakest-link pruning).
In practice $\alpha$ is selected by \emph{cross-validation} (choose the $\alpha$ with the
smallest CV error, then refit on all data), which balances bias against variance instead of
optimising training fit.
\end{solutionbox}
\fi

% ============================================================
\problem{4: Ensembles --- bagging, random forests, boosting}{16}

\begin{enumerate}[label=(\alph*),leftmargin=7mm,itemsep=2.5mm]
  \item \pp{6} Let $\hat{f}_1(x),\dots,\hat{f}_B(x)$ be $B$ identically distributed (but not
    independent) tree predictions at a fixed point $x$, each with variance $\sigma^2$ and with
    pairwise correlation $\rho$. The variance of the bagged average is
    \[
      \operatorname{Var}\Bigl(\tfrac{1}{B}\sum_{b=1}^{B}\hat{f}_b(x)\Bigr)
      =\rho\,\sigma^{2}+\frac{1-\rho}{B}\,\sigma^{2}.
    \]
    (i) Explain briefly why this formula shows that bagging reduces variance, and which term
    limits the reduction. (ii) Evaluate the variance for $B=100$, $\rho=0.6$ and $\sigma^2=4$,
    and compare it with the variance of a single tree. (iii) What value does the variance
    approach as $B\to\infty$?
  \item \pp{4} How do \emph{random forests} modify bagging at each split, and what is the
    recommended subset size $m$ for classification as a function of $p$ (evaluate for $p=25$)?
    Explain, with reference to the formula in (a), why this modification improves on bagging.
  \item \pp{4} Name the three tuning parameters of \emph{boosting} for regression trees and
    state in one sentence each what they control. What is the ``slow learning'' idea behind
    boosting?
  \item \pp{2} What does the \emph{out-of-bag} (OOB) error of a bagged model estimate, and why
    is it available essentially for free?
\end{enumerate}

\ifdefined\withsolutions
\begin{solutionbox}
\textbf{(a) Variance of the bagged average.}
(i) The second term $\frac{1-\rho}{B}\sigma^2$ is driven to $0$ by averaging over many trees
($B$ large), so averaging removes the ``independent part'' of the trees' variability ---
this is the variance reduction of bagging. The first term $\rho\sigma^2$ does \emph{not}
depend on $B$: because the bootstrap trees are grown on overlapping samples of the same data
(and share the same strong predictors), they are positively correlated, and this correlated
part of the variance cannot be averaged away. It limits the reduction.
(ii) For $B=100$, $\rho=0.6$, $\sigma^2=4$:
\[
  \operatorname{Var}=0.6\cdot 4+\frac{1-0.6}{100}\cdot 4
  =2.4+\frac{0.4\cdot 4}{100}=2.4+0.016=\mathbf{2.416},
\]
compared with $\sigma^2=4$ for a single tree --- a reduction of about $40\%$.
(iii) As $B\to\infty$ the variance approaches the floor $\rho\sigma^2=0.6\cdot 4=\mathbf{2.4}$.
(Increasing $B$ further does not overfit; it only stabilises the average.)

\medskip
\textbf{(b) Random forests.}
At each split, only a \emph{random subset} of $m$ of the $p$ predictors is considered as
splitting candidates; a fresh subset is drawn at every split. Recommended for classification:
$m\approx\sqrt{p}$, so for $p=25$: $m=\sqrt{25}=\mathbf{5}$.
Why it helps: with bagging, a dominant predictor is used in the top split of almost every tree,
so the trees look alike and $\rho$ is high. Excluding it from $\approx (p-m)/p$ of the split
searches \emph{decorrelates} the trees, i.e.\ lowers $\rho$ --- and by (a) this lowers the
irreducible floor $\rho\sigma^2$ of the ensemble variance, which mere averaging cannot touch.

\medskip
\textbf{(c) Boosting.}
\begin{itemize}[nosep,leftmargin=5mm]
  \item \emph{Number of trees $B$:} how many small trees are added sequentially; unlike
    bagging, boosting \emph{can} overfit if $B$ is too large, so $B$ is chosen by
    cross-validation.
  \item \emph{Shrinkage $\lambda$} (learning rate, e.g.\ $0.01$ or $0.001$): scales the
    contribution of each new tree and thus the speed of learning; smaller $\lambda$ usually
    needs larger $B$.
  \item \emph{Number of splits $d$ per tree} (interaction depth): the complexity of each tree
    and the interaction order of the model; often even $d=1$ (stumps, additive model) works
    well.
\end{itemize}
\emph{Slow learning:} each tree is fitted to the current \emph{residuals}, and only a small
fraction $\lambda$ of it is added to the model; the model thus improves slowly and exactly in
the regions where it still performs poorly, instead of fitting the data hard in one step ---
which tends to generalise better.

\medskip
\textbf{(d) OOB error.}
Each bootstrap sample leaves out on average about $1/3$ of the observations
(``out-of-bag''). Predicting each observation only with the trees that did \emph{not} use it,
and aggregating these predictions, yields the OOB error --- a valid estimate of the
\emph{test error} of the bagged model. It is free because it reuses the bootstrap side
product; no separate validation set or cross-validation loop is required.
\end{solutionbox}
\fi

% ============================================================
\problem{5: Neural networks by hand}{20}

Consider a tiny feed-forward network for a quantitative response with input
$x=(x_1,x_2)=(1,2)$, one hidden layer with two ReLU units
($g(z)=\max\{0,z\}$) and a linear output unit:
\[
  A_1=g\bigl(-1+1\,x_1+2\,x_2\bigr),\qquad
  A_2=g\bigl(-3-2\,x_1+1\,x_2\bigr),\qquad
  f(x)=1+2\,A_1+4\,A_2 .
\]

\begin{enumerate}[label=(\alph*),leftmargin=7mm,itemsep=2.5mm]
  \item \pp{8} Perform the forward pass step by step: compute both pre-activations, both
    activations $A_1,A_2$ and the network output $f(x)$. Comment briefly on what the ReLU did
    to the second unit and why such non-linear activations are essential.
  \item \pp{6} A classification network with $K=3$ classes produces the logits (output-layer
    pre-activations) $z=(z_1,z_2,z_3)=(1.0,\ 0.0,\ -1.0)$ for one observation whose true class
    is class~1. Compute all three softmax probabilities and the cross-entropy loss
    $-\ln \hat{p}_{\text{true}}$.
    \emph{(Use $e^{1}=2.718$, $e^{0}=1.000$, $e^{-1}=0.368$ and $\ln 4.086=1.408$.)}
  \item \pp{4} Count the parameters of a fully connected network with $p=4$ inputs, one hidden
    layer with $k=3$ units and an output layer with $K=2$ units, where every hidden and output
    unit has a bias. Show the count separately for the two weight layers.
  \item \pp{2} A convolutional layer receives a $32\times 32$ image and uses filters of size
    $F=5$ with padding $P=2$ and stride $S=1$. Compute the spatial size of the output feature
    map from $\bigl(W-F+2P\bigr)/S+1$.
\end{enumerate}

\ifdefined\withsolutions
\begin{solutionbox}
\textbf{(a) Forward pass.}
Pre-activations:
\[
  z_1=-1+1\cdot 1+2\cdot 2=-1+1+4=4,\qquad
  z_2=-3-2\cdot 1+1\cdot 2=-3-2+2=-3 .
\]
ReLU activations:
\[
  A_1=g(4)=\max\{0,4\}=4,\qquad A_2=g(-3)=\max\{0,-3\}=0 .
\]
Output:
\[
  f(x)=1+2\cdot 4+4\cdot 0=1+8+0=\mathbf{9}.
\]
The ReLU \emph{switched the second unit off}: its negative pre-activation is clipped to $0$,
so unit 2 contributes nothing for this input (but would for other inputs). This thresholding
is the source of non-linearity: without a non-linear $g$, the network would collapse into a
single linear function of $x$, no matter how many layers it has.

\medskip
\textbf{(b) Softmax and cross-entropy.}
Exponentials: $e^{1.0}=2.718$, $e^{0.0}=1.000$, $e^{-1.0}=0.368$; their sum is
$2.718+1.000+0.368=4.086$. Softmax probabilities:
\[
  \hat{p}_1=\frac{2.718}{4.086}=0.665,\qquad
  \hat{p}_2=\frac{1.000}{4.086}=0.245,\qquad
  \hat{p}_3=\frac{0.368}{4.086}=0.090,
\]
which sum to $1.000$. Cross-entropy loss for true class 1:
\[
  -\ln\hat{p}_1=-\ln\frac{e^{1}}{4.086}=-\bigl(1-\ln 4.086\bigr)=1.408-1=\mathbf{0.408}.
\]
(Equivalently $-\ln 0.665\approx 0.408$; a perfect prediction $\hat{p}_1=1$ would give loss 0.)

\medskip
\textbf{(c) Parameter count for $4\to 3\to 2$.}
Input $\to$ hidden: $4\cdot 3=12$ weights $+\,3$ biases $=15$.
Hidden $\to$ output: $3\cdot 2=6$ weights $+\,2$ biases $=8$.
Total: $15+8=\mathbf{23}$ parameters.

\medskip
\textbf{(d) Convolution output size.}
\[
  \frac{W-F+2P}{S}+1=\frac{32-5+2\cdot 2}{1}+1=\frac{31}{1}+1=32,
\]
so the output feature map is $\mathbf{32\times 32}$ --- padding $P=2$ with $F=5$, $S=1$ is
exactly the ``same'' convolution that preserves the spatial size.
\end{solutionbox}
\fi

% ============================================================
\problem{6: Multiple testing}{20}

\begin{enumerate}[label=(\alph*),leftmargin=7mm,itemsep=2.5mm]
  \item \pp{4} A lab performs $m=20$ \emph{independent} hypothesis tests, each at level
    $\alpha=0.05$, and suppose all 20 null hypotheses are true. Define the family-wise error
    rate (FWER) and compute it for this situation
    \emph{(use $(0.95)^{20}=0.3585$)}. Interpret the result in one sentence.
  \item \pp{7} A research group screens $m=6$ candidate biomarkers; the ordered $p$-values are
    \begin{center}
    \begin{tabular}{lcccccc}
    \toprule
    $j$ & 1 & 2 & 3 & 4 & 5 & 6\\
    \midrule
    $p_{(j)}$ & 0.001 & 0.004 & 0.011 & 0.030 & 0.045 & 0.230\\
    \bottomrule
    \end{tabular}
    \end{center}
    Control the FWER at $\alpha=0.05$ (i) with \emph{Bonferroni} and (ii) with the
    \emph{Holm step-down} procedure. Give the relevant thresholds, state for each method
    exactly which hypotheses are rejected, and explain why Holm can never reject fewer
    hypotheses than Bonferroni.
  \item \pp{6} Apply the \emph{Benjamini--Hochberg} (BH) procedure with FDR level $q=0.05$ to
    the same six $p$-values: give the thresholds $jq/m$, determine which hypotheses are
    rejected, and state precisely what ``FDR control at $q=0.05$'' guarantees for the
    resulting set of rejections.
  \item \pp{3} A genomics group screens $20{,}000$ genes in an \emph{exploratory} study; the
    most promising hits will be re-tested in an independent follow-up experiment. Should they
    control the FWER or the FDR? Justify your recommendation.
\end{enumerate}

\ifdefined\withsolutions
\begin{solutionbox}
\textbf{(a) FWER for $m=20$ independent tests.}
$\text{FWER}=\Pr(\text{at least one Type I error})=\Pr(V\ge 1)$, where $V$ is the number of
true nulls that are (falsely) rejected. With independence and all nulls true:
\[
  \text{FWER}=1-\Pr(\text{no false rejection})=1-(1-0.05)^{20}=1-(0.95)^{20}
  =1-0.3585=\mathbf{0.6415}.
\]
Interpretation: even though each single test is performed at the $5\%$ level, there is a
$64\%$ chance of at least one false discovery among the 20 tests --- testing many hypotheses
without adjustment almost guarantees spurious findings.

\medskip
\textbf{(b) Bonferroni and Holm at $\alpha=0.05$.}
\emph{Bonferroni:} reject $H_{(j)}$ iff $p_{(j)}\le\alpha/m=0.05/6=0.00833$.
$p_{(1)}=0.001\le 0.00833$ \checkmark, $p_{(2)}=0.004\le 0.00833$ \checkmark,
$p_{(3)}=0.011>0.00833$ $\times$.
$\Rightarrow$ Bonferroni rejects $H_{(1)}$ and $H_{(2)}$: \textbf{2 rejections}.

\emph{Holm:} compare $p_{(j)}$ with $\alpha/(m-j+1)$ step by step and stop at the first failure:
\begin{center}
\begin{tabular}{lcccccc}
\toprule
$j$ & 1 & 2 & 3 & 4 & 5 & 6\\
\midrule
$p_{(j)}$ & 0.001 & 0.004 & 0.011 & 0.030 & 0.045 & 0.230\\
$\alpha/(m-j+1)$ & 0.00833 & 0.0100 & 0.0125 & 0.0167 & 0.0250 & 0.0500\\
reject? & \checkmark & \checkmark & \checkmark & stop & --- & ---\\
\bottomrule
\end{tabular}
\end{center}
$0.001\le 0.00833$, $0.004\le 0.0100$, $0.011\le 0.0125$, but $0.030>0.0167$: stop.
$\Rightarrow$ Holm rejects $H_{(1)},H_{(2)},H_{(3)}$: \textbf{3 rejections}.

\emph{Why Holm $\supseteq$ Bonferroni:} Holm's thresholds $\alpha/(m-j+1)$ are $\ge\alpha/m$
for every $j$ (with equality only at $j=1$), so every $p$-value that passes Bonferroni also
passes Holm's less stringent later hurdles --- Holm is uniformly more powerful, yet still
controls the FWER at $\alpha$ without any independence assumption.

\medskip
\textbf{(c) Benjamini--Hochberg at $q=0.05$.}
Thresholds $jq/m=j\cdot 0.05/6$:
\begin{center}
\begin{tabular}{lcccccc}
\toprule
$j$ & 1 & 2 & 3 & 4 & 5 & 6\\
\midrule
$p_{(j)}$ & 0.001 & 0.004 & 0.011 & 0.030 & 0.045 & 0.230\\
$jq/m$ & 0.00833 & 0.0167 & 0.0250 & 0.0333 & 0.0417 & 0.0500\\
$p_{(j)}\le jq/m$? & yes & yes & yes & \textbf{yes} & no & no\\
\bottomrule
\end{tabular}
\end{center}
$L=$ largest $j$ with $p_{(j)}\le jq/m$: $j=4$ (since $0.030\le 0.0333$, while
$0.045>0.0417$ and $0.230>0.0500$). BH rejects \emph{all} $H_{(j)}$ with $j\le L$:
$H_{(1)},\dots,H_{(4)}$: \textbf{4 rejections}.
Guarantee: $\text{FDR}=E\bigl[V/\max(R,1)\bigr]\le q=0.05$ --- \emph{on average} at most
$5\%$ of the rejected hypotheses are false discoveries. It does \emph{not} guarantee that at
most $5\%$ of these particular 4 rejections are false, nor does it bound the probability of
making any false rejection (that would be FWER control).

\medskip
\textbf{(d) Exploratory genomics screening.}
Control the \textbf{FDR}. With $m=20{,}000$ tests, FWER control (e.g.\ Bonferroni at
$0.05/20{,}000=2.5\cdot 10^{-6}$) guards against even a single false positive and therefore
has extremely low power --- almost no true signals would survive. In an exploratory screen
the cost of a few false leads is low, because every hit is validated in an independent
follow-up experiment; what matters is that the \emph{list} of discoveries is mostly reliable.
FDR control at, say, $q=0.05$--$0.2$ delivers exactly that: a rich list of candidates of
which on average at most a known fraction is false.
\end{solutionbox}
\fi

% ============================================================
\problem{7: Cumulative essentials (Chapters 2--6)}{20}

\begin{enumerate}[label=(\alph*),leftmargin=7mm,itemsep=2.5mm]
  \item \pp{4} True or false --- justify briefly: ``A more flexible statistical learning method
    always yields a lower \emph{test} error than a less flexible one.''
  \item \pp{4} A logistic regression for loan default,
    $\log\bigl(\frac{p(x)}{1-p(x)}\bigr)=\beta_0+\beta_1 x$ with $x=$ number of missed
    payments, yields $\hat{\beta}_1=0.7$. Interpret the effect of \emph{one additional missed
    payment} on the \emph{odds} of default \emph{(use $e^{0.7}=2.014$)}. Why can the effect on
    the \emph{probability} of default not be stated as a single number?
  \item \pp{4} Give one reason to prefer $k$-fold cross-validation with $k=5$ or $10$ over
    \emph{LOOCV}, and one reason to prefer it over a \emph{single validation split}.
  \item \pp{2} You have $p=2{,}000$ predictors but expect only a handful to be truly related to
    the response. Ridge or lasso --- which penalty is more appropriate, and why?
  \item \pp{6} Match each scenario to exactly one method from
    \{\emph{linear regression}, \emph{lasso}, \emph{logistic regression},
    \emph{random forest}\} (each method used exactly once) and justify each choice in one
    sentence:
    \begin{enumerate}[label=(\roman*),nosep,topsep=1mm]
      \item Estimate how much an additional square metre adds to the price of an apartment,
        with a confidence interval, from $n=2{,}000$ sales.
      \item Predict a continuous drug response from $p=5{,}000$ gene-expression measurements
        for $n=100$ patients; clinicians want a short list of relevant genes.
      \item Predict whether a loan applicant will default (yes/no) and report to the regulator
        how each applicant characteristic changes the odds of default.
      \item Predict machine failures from hundreds of sensor variables with strong non-linear
        effects and interactions; only predictive accuracy matters.
    \end{enumerate}
\end{enumerate}

\ifdefined\withsolutions
\begin{solutionbox}
\textbf{(a) False.}
Expected test error $=$ variance $+$ bias$^2$ $+$ irreducible error. More flexibility reduces
bias but \emph{increases} variance; once the variance increase outweighs the bias reduction,
the test error rises again --- the characteristic U-shape. So a more flexible method wins only
up to the sweet spot; e.g.\ when the true $f$ is nearly linear or $n$ is small, a flexible
method overfits and a simple linear fit has \emph{lower} test error. (Only the \emph{training}
error always decreases with flexibility.)

\medskip
\textbf{(b) Odds ratio.}
One additional missed payment increases the log-odds by $\hat{\beta}_1=0.7$, i.e.\ it
\emph{multiplies the odds} of default by $e^{0.7}=2.014$ --- the odds roughly double,
regardless of the starting value of $x$.
The effect on the \emph{probability} depends on where the observation sits on the S-shaped
logistic curve: the same change in log-odds shifts $p(x)$ a lot near $p\approx 0.5$ but very
little when $p$ is close to $0$ or $1$ --- hence no single number describes it.

\medskip
\textbf{(c) $k$-fold CV with $k=5$ or $10$.}
\emph{vs.\ LOOCV:} the $n$ LOOCV fits are trained on almost identical data sets, so their
errors are highly positively correlated; averaging strongly correlated quantities reduces
variance little, so the LOOCV estimate has \emph{higher variance} than 5- or 10-fold CV
(it is also computationally expensive: $n$ fits instead of $k$, absent shortcut formulas).
\emph{vs.\ a single validation split:} the validation estimate depends strongly on the random
split and overestimates the test error because only part of the data is used for training
(\emph{bias}); $k$-fold CV trains on $(k-1)/k$ of the data and averages over $k$ folds, giving
a less biased and more stable estimate.

\medskip
\textbf{(d) Lasso.}
The $\ell_1$ penalty sets many coefficients \emph{exactly to zero} and thus performs variable
selection --- ideal when the truth is sparse (a handful of signals among 2{,}000 predictors),
and it yields an interpretable small model. Ridge shrinks coefficients but keeps all
$p=2{,}000$ predictors in the model and tends to lose when most true coefficients are zero.

\medskip
\textbf{(e) Matching.}
\begin{itemize}[nosep,leftmargin=5mm]
  \item (i) \textbf{Linear regression}: quantitative response, goal is \emph{inference}
    (effect size with confidence interval) --- exactly what OLS with standard errors provides.
  \item (ii) \textbf{Lasso}: $p\gg n$ makes OLS infeasible; the $\ell_1$ penalty regularises
    and selects a short list of genes.
  \item (iii) \textbf{Logistic regression}: binary response with interpretable coefficients
    ($e^{\hat\beta_j}$ = odds ratios) --- suitable for regulatory reporting.
  \item (iv) \textbf{Random forest}: captures non-linearities and interactions automatically,
    is accurate and robust with many mixed predictors; interpretability is not required.
\end{itemize}
\end{solutionbox}
\fi

\vspace{6mm}
\begin{center}
\rule{0.35\textwidth}{0.4pt}\\[1mm]
\textit{End of the examination --- good luck!}
\end{center}

\end{document}
